\documentclass[11pt]{amsart}

\usepackage[margin=1.08in]{geometry}
\usepackage{amsmath,amssymb,amsthm,mathtools}
\usepackage{booktabs,array,longtable,xcolor}
\usepackage{microtype}
\usepackage{aliascnt}
\usepackage[hidelinks]{hyperref}
\usepackage[nameinlink,noabbrev]{cleveref}

\newtheorem{theorem}{Theorem}[section]
\newaliascnt{proposition}{theorem}
\newtheorem{proposition}[proposition]{Proposition}
\aliascntresetthe{proposition}
\newaliascnt{corollary}{theorem}
\newtheorem{corollary}[corollary]{Corollary}
\aliascntresetthe{corollary}
\newaliascnt{lemma}{theorem}
\newtheorem{lemma}[lemma]{Lemma}
\aliascntresetthe{lemma}
\newaliascnt{conjecture}{theorem}
\newtheorem{conjecture}[conjecture]{Conjecture}
\aliascntresetthe{conjecture}
\theoremstyle{definition}
\newaliascnt{definition}{theorem}
\newtheorem{definition}[definition]{Definition}
\aliascntresetthe{definition}
\newaliascnt{computation}{theorem}
\newtheorem{computation}[computation]{Exact computation}
\aliascntresetthe{computation}
\theoremstyle{remark}
\newaliascnt{remark}{theorem}
\newtheorem{remark}[remark]{Remark}
\aliascntresetthe{remark}

\crefname{computation}{Exact computation}{Exact computations}
\Crefname{computation}{Exact computation}{Exact computations}

\newcommand{\QQ}{\mathbf Q}
\newcommand{\ZZ}{\mathbf Z}
\newcommand{\FF}{\mathbf F}
\newcommand{\PP}{\mathbf P}
\newcommand{\Gal}{\operatorname{Gal}}
\newcommand{\PSL}{\operatorname{PSL}}
\newcommand{\Spl}{\operatorname{Spl}}
\newcommand{\Sym}{\operatorname{Sym}}
\newcommand{\Fix}{\operatorname{Fix}}
\newcommand{\Res}{\operatorname{Res}}
\newcommand{\disc}{\operatorname{disc}}
\newcommand{\ord}{\operatorname{ord}}
\newcommand{\Out}{\operatorname{Out}}
\newcommand{\Aut}{\operatorname{Aut}}
\newcommand{\Inn}{\operatorname{Inn}}
\newcommand{\Span}{\operatorname{span}}
\newcommand{\M}{M_{23}}
\newcommand{\Mtt}{M_{23}}
\newcommand{\GQ}{\mathfrak G_{\QQ}}
\newcommand{\Cx}{C_{M_{23}}(x)}

\title[Minimal projection, Fano arithmetic, and the Hurwitz scheme]
  {Minimal-degree projection, Fano-plane arithmetic, and the finite
   Hurwitz scheme of a regular Mathieu $M_{23}$-cover}
\author{François Le Lay}
\address{Applied Science Team, Reflexivity, 589 5th Avenue, Suite 606A,
New York, NY 10017}
\email{francois@reflexivity.com}
\urladdr{https://www.reflexivity.com}
\thanks{Copyright 2026 Knabble, Inc. (dba Reflexivity). Licensed under the
Creative Commons Attribution-NonCommercial-ShareAlike 4.0 International
License (\url{https://creativecommons.org/licenses/by-nc-sa/4.0/}).}
\date{}

\hypersetup{
  pdftitle={Minimal-degree projection, Fano-plane arithmetic, and the finite Hurwitz scheme of a regular Mathieu M23-cover},
  pdfauthor={François Le Lay},
  pdfsubject={Minimal-degree projection, branch-fibre arithmetic, and the finite Hurwitz scheme of a regular M23-cover}
}

\subjclass[2020]{Primary 12F12; Secondary 11G32, 14H30, 20B25, 11R32}
\keywords{Inverse Galois problem, Mathieu group $M_{23}$, Fano plane,
minimal-degree projection, branch fibre, finite Hurwitz scheme,
arithmetically equivalent fields}

\begin{document}

\begin{abstract}
Starting from the regular degree-$23$ $M_{23}$-cover constructed by Huang,
Jackson, Lee, Poonen, Pries, and Zhang, we give one exact account of three
structures surrounding it: an explicit realization of the minimal-degree
projection of its source curve, the arithmetic of the branch fibre over
$T=0$, and the finite inner Hurwitz scheme parametrizing normalized
three-point $M_{23}$-covers with branch-cycle classes $(2A,23A,23B)$.

\smallskip
Huang et al. prove that the genus-$4$ source curve has $\QQ$-gonality $4$,
but leave a degree-$4$ coordinate uncomputed.  We make their anticipated
construction explicit: two canonical differentials have a ratio $W$ of
degree $4$, and elimination gives a primitive equation
\[
 P(T,W)\in\ZZ[T,W],\qquad \deg_WP=23,\quad \deg_TP=4.
\]
An exact function-field identity, together with irreducibility modulo $31$,
proves that this equation defines the original cover.

\smallskip
Over the rational point $T=0$, the integral source equation factors as
$23^4H_7R_8^2$, with irreducible factors of degrees $7$ and $8$.  The roots
of $H_7$ carry the point set of a Fano plane; its lines yield a
nonisomorphic arithmetically equivalent septic field.  The roots of $R_8$
form an affine $\FF_2^3$-torsor, and separating their two local branches
completes a Galois tower with group $2^4{:}\PSL_3(2)$.  The change of
generator $W=J_0(T,V)/J_1(T,V)$ is regular throughout this fibre and
\[
 P(0,W)=-23^3\widetilde H_7(W)\widetilde R_8(W)^2.
\]
It therefore preserves the Fano plane, affine cube, and local field tower
exactly.

\smallskip
The Fano-plane structure occurs in a single branch fibre of the distinguished
cover.  Separately, over $K_0=\QQ(\sqrt{-23})$, exact reconstruction of the
finite Hurwitz scheme gives the coordinate algebra $K_0\times L$ with
$[L:K_0]=6$, together with the seven canonical genus-$4$ curves and their
degree-$23$ maps.  Certified branch cycles show that the maps exhaust the
seven Nielsen classes and all have monodromy $M_{23}$.  The relative Galois
closure of $L/K_0$ has group $S_6$, acting naturally on the six conjugate
maps.  Together these calculations separate the arithmetic intrinsic to the
chosen branch fibre from the Galois action on the Hurwitz scheme; a conceptual
explanation of the known fixed Nielsen class remains open.

\smallskip
Finally, for every integer $t\equiv2\pmod{319}$, the primitive part of
$P(t,W)$ has Galois group $M_{23}$; controlled new ramification selects an
infinite mutually independent subfamily inside this explicit progression.
Exact SageMath, Singular, Magma, GAP, Arb, and PARI/GP certificates accompany
the computations.
\end{abstract}

\maketitle

\clearpage
\setcounter{tocdepth}{1}
\tableofcontents

\section{Introduction}

Let $X\to\PP^1_T$ be the degree-$23$ cover constructed in
\cite{HJLPPZ}.  Its splitting field over $\QQ(T)$ is regular with Galois
group $\M$, and its source curve has genus $4$.  The three branch points are
\[
 0,\qquad \sqrt{-23},\qquad -\sqrt{-23},
\]
with local monodromy classes $(2A,23A,23B)$.  In particular the inertia
involution at $T=0$ has cycle shape $1^7 2^8$: seven sheets remain fixed
and eight pairs collide.

\smallskip
Two distinct seven-element Galois sets enter the argument.  The seven
unramified points in the fibre over $T=0$ carry the Fano-plane incidence
studied below.  The seven geometric points of the inner Hurwitz scheme
parametrize the covers in the Nielsen-class problem.  The first records
arithmetic within the distinguished cover, whereas the second records
arithmetic among the seven covers.

\smallskip
Here and throughout, the phrase \emph{fibre above the rational branch
point} refers to the fibre of $X\to\PP^1_T$ above
$0\in\PP^1(\QQ)$; it does not assert that the points of that fibre are
$\QQ$-rational.  This is a characteristic-zero branch fibre.  We reserve
\emph{closed fibre} (often called a special fibre) for a model over a local
base, and we refer to passage modulo a prime as reduction.

\smallskip
The paper follows one cover from its global geometry to one structured
branch fibre, then returns to the moduli problem containing all seven covers.
A uniform family of arithmetic specializations is a separate application of
the same minimal equation.  Write $F(T,V)\in\ZZ[T,V]$ for the integral source
equation.  Huang et al. prove that $X$ admits a $\QQ$-defined map of degree
$4$, that no lower-degree map is defined over $\QQ$, and remark that one could
replace their degree-$8$ coordinate by such a function
\cite[Remark~3.4]{HJLPPZ}.  They do not compute it.  Our first result carries
out that anticipated construction and certifies the resulting equation
exactly.

\begin{theorem}[Explicit realization of the minimal-degree projection]
\label{thm:main}
There are explicit polynomials $J_0,J_1\in\ZZ[T,V]$ and a primitive
irreducible polynomial $P\in\ZZ[T,W]$ such that
\[
 \deg_WP=23,\qquad \deg_TP=4,
\]
and, in $\QQ(X)$,
\[
 W=\frac{J_0(T,V)}{J_1(T,V)},\qquad P(T,W)=0,
 \qquad \QQ(T,W)=\QQ(T,V).
\]
Consequently $P$ defines the same regular $M_{23}$-cover.  The projection
$X\to\PP^1_W$ has degree $4$, and therefore explicitly realizes the minimum
over $\QQ$ established in \cite[Remark~3.4]{HJLPPZ}.
\end{theorem}

The fibre above the rational branch point contains the local arithmetic of
the inertia centralizer.  In the original coordinate its defining
polynomial factors as
$23^4H_7R_8^2$: the seven simple roots support a Fano plane and the eight
double roots support an affine cube.  The following theorem connects that
local structure directly to the minimal equation.

\begin{theorem}[Transport of the fibre above $T=0$]\label{thm:branch-fibre-transport}
There are primitive irreducible polynomials
$\widetilde H_7,\widetilde R_8\in\ZZ[W]$ of degrees $7$ and $8$ such that
\[
\begin{aligned}
 F(0,V)&=23^4H_7(V)R_8(V)^2,\\
 P(0,W)&=-23^3\widetilde H_7(W)\widetilde R_8(W)^2.
\end{aligned}
\]
Moreover
\[
 \gcd\bigl(J_1(0,V),H_7(V)R_8(V)\bigr)=1,
\]
and substitution $W=J_0(0,V)/J_1(0,V)$ induces isomorphisms
\[
\begin{aligned}
 \QQ[W]/(\widetilde H_7)&\xrightarrow{\ \sim\ }\QQ[V]/(H_7),\\
 \QQ[W]/(\widetilde R_8)&\xrightarrow{\ \sim\ }\QQ[V]/(R_8).
\end{aligned}
\]
These isomorphisms give Galois-equivariant bijections between the respective
sets of geometric roots.
\end{theorem}

Huang et al. numerically compute all seven complex covers and identify the
Galois-fixed one.  The branch-fibre calculation here concerns only that
distinguished cover.  To determine the arithmetic structure and Galois action
on the full seven-cover moduli problem, we reconstruct its finite étale
$K_0$-scheme exactly.

\begin{theorem}[Finite inner Hurwitz scheme]\label{thm:hurwitz-summary}
Put $K_0=\QQ(\sqrt{-23})$, and let $\mathcal H_{\mathrm{in}}$ be the reduced
inner Hurwitz scheme of normalized three-point covers with branch-cycle
classes $(2A,23A,23B)$.  There is a degree-$6$ field extension $L/K_0$ and
an isomorphism
\[
 \mathcal H_{\mathrm{in}}\simeq\operatorname{Spec}(K_0\times L).
\]
The relative Galois closure of $L/K_0$ has group $S_6$.  Over this coordinate
algebra there are seven explicit canonical genus-$4$ curves with exact
degree-$23$ maps; their branch cycles exhaust the seven Nielsen classes and
each map has monodromy $M_{23}$.
\end{theorem}

Finally, the minimal equation has a uniform arithmetic consequence away
from the branch locus.

\begin{theorem}[Uniform specialization]\label{thm:specialization}
For every integer
\[
 t\equiv2\pmod{319},
\]
the primitive part of $P(t,W)$ has degree $23$, and its splitting field over
$\QQ$ has Galois group $M_{23}$ in its natural degree-$23$ action.
\end{theorem}

The argument has three principal stages and one arithmetic application.
\begin{enumerate}
\item \emph{Explicit minimal projection.}  Starting from the degree-$4$
  existence and minimality result of Huang et al., adjoint differentials on
  the $84$-node good reduction of their plane model produce an explicit
  canonical pencil and function $W$.
\item \emph{Branch fibre.}  At $T=0$, the $23$ sheets become seven fixed sheets
  and eight colliding pairs.  Their residue arithmetic is the Fano plane and
  the affine cube; \cref{thm:branch-fibre-transport} transports both
  through $W$.
\item \emph{Finite Hurwitz scheme.}  Exact reconstruction, elimination, and
  certified root continuation recover all seven covers and the full
  $1+6$ arithmetic decomposition of their moduli scheme.
\item \emph{Arithmetic specialization.}  Moving away from the branch fibre in the
  congruence class $t\equiv2\pmod{319}$ recovers the full group $M_{23}$ for
  every integer specialization.
\end{enumerate}

The coordinate comparison transports the Fano and affine structures from the
original plane equation to the minimal-degree coordinate.  Independently, the
Hurwitz-scheme calculation determines the actual $1+6$ Galois decomposition.
A conceptual construction of its degree-one component from the branch-fibre
data remains open.

\subsection{Prior work and contributions}
\label{sec:prior-work}

The regular $M_{23}$-cover, its source equation and branch cycles, the
numerical computation of the seven complex covers, and the identification of
the Galois-fixed Nielsen class are due to Huang et al.\ \cite{HJLPPZ}.  They
also prove that the source curve has $\QQ$-gonality $4$ and observe that a
degree-$4$ coordinate could replace their degree-$8$ coordinate, without
computing such a function.  The branch-fibre and Hurwitz-scheme calculations
in this paper begin from their cover and its accompanying data
\cite{M23data}.

The braid-orbit calculation behind the seven-element Nielsen class goes back
to H\"afner \cite{BraidOrbits}.  The numerical three-point-cover method used
by Huang et al. is that of Klug--Musty--Schiavone--Voight \cite{KMSV},
implemented in the Belyi-map database of
Musty--Schiavone--Sijsling--Voight \cite{BelyiDB}.  The point--line Gassmann
pair and the flag description of transvections in $\PSL_3(2)$ are standard.
Likewise, the abstract solvability of the split embedding problem with
abelian kernel is not at issue here \cite{LegrandEmbedding}; our calculation
identifies the fields cut out by this particular branch fibre.  Huang et al.
already obtain abstractly an infinite mutually independent family of
$M_{23}$-extensions \cite[Corollary~1.4(d)]{HJLPPZ}, and general
Hilbert--Grunwald methods provide arithmetic progressions with prescribed
specialization behavior \cite{DebesGhazi,Legrand}.

Relative to this prior work, the contributions here are the explicit
minimal-degree equation and ruling field; the exact arithmetic and
obstruction theory of the distinguished branch fibre; the exact
$K_0\times L$ Hurwitz algebra, its $S_6$ Galois closure, and the seven exact
maps; and the all-members congruence class modulo $319$ with controlled
ramification.  The detailed attribution, limitations, and data provenance
are recorded in \cref{sec:scope}.

\section{Explicitly constructing the minimal-degree projection}
\label{sec:adjoints}

Set $D=T^2+23$.  The coefficient of $V^{23}$ in $F$ is $D^4$, so
\[
 \Phi:=F,\qquad \widehat F:=D^{-4}F\in\QQ(T)[V]
\]
defines the same function field and is monic of degree $23$ in $V$.  We
write $C_0$ for the affine plane curve $\Phi=0$ and $X$ for its
normalization.  The compactification under consideration is the closure in
$\PP^1_T\times\PP^1_V$.

\subsection{The two points above \texorpdfstring{$T^2+23=0$}{T2+23=0}}

Let $b,c\in X(\overline{\QQ})$ be the unique points over the roots of $D$.
They are conjugate over $\QQ$ and each has ramification index $23$ for the
$T$-map \cite[Section~3]{HJLPPZ}.

\begin{lemma}[Valuations at the two ramification points]\label{lem:valuations}
At either $q\in\{b,c\}$ one has
\[
 \ord_q(T-T(q))=23,\qquad \ord_q(D)=23,
 \qquad \ord_q(V)=-4,
\]
and
\[
 \ord_q\!\left(\frac{dT}{\Phi_V}\right)=18.
\]
\end{lemma}

\begin{proof}
The first two equalities follow from total ramification and the fact that
$D$ has a simple zero at $T(q)$.  The coefficients of $\Phi$, viewed as a
polynomial in $V$, have a Newton-polygon edge joining $(0,0)$ and $(23,4)$ at
$D=0$.  Hence its unique branch has $\ord_q(V)=-4$.

For completeness, the terms of $\Phi_V$ on this branch have orders
\[
 23\ord_D(i a_i)-4(i-1),
\]
where $a_i(T)$ is the coefficient of $V^i$.  The minimum is $4$, attained
uniquely by differentiating the $D^4V^{23}$ term; all other orders are at
least $5$.  Thus $\ord_q(\Phi_V)=4$.  Finally
$\ord_q(dT)=23-1=22$, so $\ord_q(dT/\Phi_V)=22-4=18$.
\end{proof}

The calculation is reproduced directly from the normalized coefficient
table by \path{verification/verify_boundary_valuations.gp}.

For a polynomial $A(T,V)$ consider the adjoint differential
\[
 \omega_A=\frac{A(T,V)\,dT}{\Phi_V(T,V)}.
\]
The interior of the bidegree-$(8,23)$ rectangle gives the usual support
$0\leq i\leq6$, $0\leq j\leq21$ for monomials $T^iV^j$.  Regularity at
$b$ and $c$ sharpens this $7\cdot22=154$-dimensional space substantially.

\begin{lemma}[Adjoints regular at $b$ and $c$]\label{lem:adjoint-space}
The differentials $\omega_A$ having interior support and regular at both
$b$ and $c$ are represented by the $88$-dimensional space
\begin{align*}
\mathcal A={}&
 \bigoplus_{j=0}^{4} V^j\QQ[T]_{\leq6}
 \;\oplus\!
 \bigoplus_{j=5}^{10} DV^j\QQ[T]_{\leq4}\\
&\oplus\!
 \bigoplus_{j=11}^{16}D^2V^j\QQ[T]_{\leq2}
 \;\oplus\!
 \bigoplus_{j=17}^{21}D^3V^j\QQ.
\end{align*}
In particular,
\[
 \dim\mathcal A=35+30+18+5=88.
\]
\end{lemma}

\begin{proof}
Write $A=\sum_{j=0}^{21}a_j(T)V^j$.  By
\cref{lem:valuations}, a term $D^m c(T)V^j\,dT/\Phi_V$ has order
\[
 23m-4j+18
\]
at $b$ and $c$, unless $c$ also vanishes there.  Regularity therefore
requires
\[
 m\geq\left\lceil\frac{4j-18}{23}\right\rceil,
\]
giving respectively $m=0,1,2,3$ on the four displayed blocks.  Terms having
distinct exponents $j$ cannot cancel: their orders are distinct modulo $23$.
Divisibility by $D^m$ lowers the residual $T$-degree from $6$ to $6-2m$.
Counting coefficients gives the stated dimension.
\end{proof}

\subsection{The affine nodes and the canonical system}

At an ordinary node, Rosenlicht's adjoint condition \cite{Rosenlicht} is the
vanishing of the numerator $A$.  Thus each node contributes one linear
evaluation condition, and the rank computation below verifies independence.
We impose the conditions in one stroke using the reduced singular scheme
modulo $31$.

\begin{computation}[The node graph modulo 31]\label{comp:nodes}
Over $\FF_{31}$ the Jacobian ideal
\[
 (\Phi,\Phi_T,\Phi_V)
\]
has lexicographic basis
\[
 h_{84}(T),\qquad V-s(T),
\]
where $\deg h_{84}=84$ and $h_{84}$ is squarefree.  The determinant of the
Hessian is a unit on this scheme.  Thus the affine singular scheme consists
geometrically of $84$ reduced ordinary nodes.
\end{computation}

This statement is certified by
\path{verification/verify_nodes_mod31.sing}; the explicit $h_{84}$ and
$s$ are in \path{data/node_graph_mod31.inc}.  Evaluation of the basis in
\cref{lem:adjoint-space} on the node graph gives an $84\times88$ matrix over
$\FF_{31}$.

\begin{computation}[Canonical kernel and degree-$4$ pencil]\label{comp:kernel}
The node-evaluation matrix has rank $84$, so its kernel has dimension $4$.
This is the canonical adjoint space $H^0(X,K_X)$.  On this kernel, vanishing
at $b+c$ has rank $2$, leaving the two-dimensional space
$H^0(X,K_X-b-c)$ and hence a pencil.
\end{computation}

\begin{proof}[Explanation of the two conditions at $b$ and $c$]
Among the blocks of \cref{lem:adjoint-space}, the only term that can have
order zero at $b$ or $c$ is
\[
 D^2(c_0+c_1T+c_2T^2)V^{16}.
\]
Its vanishing at both roots of $D$ is equivalent to divisibility of the
quadratic coefficient by $D$, namely
\[
 c_1=0,\qquad c_0=23c_2.
\]
These two conditions have rank $2$ on the four-dimensional node kernel.
The ranks are reproduced by
\path{verification/verify_adjoint_mod31.py}.
\end{proof}

The successive dimensions are therefore
\[
154\longrightarrow88\longrightarrow4\longrightarrow2.
\]
The last two arrows are certified by full-rank minors modulo $31$.  This
good-reduction calculation is the modular construction step; it does not by
itself certify the characteristic-zero equation.  That role is played by
the exact function-field identity and irreducibility proof in
\cref{sec:certificate}.

If $\omega_0$ and $\omega_1$ form a basis of
$H^0(X,K_X-b-c)$, then their ratio $W=\omega_0/\omega_1$ is a rational
function on $X$.  The resulting row-reduced-echelon-form normalized pencil
was reconstructed over $\QQ$ by the Chinese remainder theorem.  The least
common multiple of the reconstructed coefficient denominators is $91$;
clearing it gives $J_0,J_1\in\ZZ[T,V]$.  Their degrees in $V$ are $20$ and
$21$, respectively, and each has degree $6$ in $T$.

\subsection{Reconstruction of the bidegree \texorpdfstring{$(23,4)$}{(23,4)} equation}
\label{sec:reconstruction}

Set $W=J_0/J_1$.  Eliminating $V$ between $\Phi(T,V)$ and
$J_0(T,V)-WJ_1(T,V)$ means extracting the desired factor from
\[
 \Res_V\!\bigl(\Phi(T,V),J_0(T,V)-WJ_1(T,V)\bigr).
\]
We performed this calculation modulo good primes in a stable
row-reduced-echelon-form (RREF) normalization, reconstructed the
coefficients, and then discarded the
numerical evidence in favor of the exact certificate below.

For transparency, the equation used a $575$-bit CRT modulus.  All $120$
coefficients reconstructed and reduced correctly at the independent holdout
prime $500015017$.  The pencil used a $604$-bit modulus and passed independent
holdout prime $137$.  These holdouts are evidence about reconstruction, not
ingredients in the proof of \cref{thm:main}; their complete records are
\path{data/f4_crt_reconstruction.json} and
\path{data/pencil_crt_reconstruction.json}.  As an auxiliary geometric check,
the reduction of $P$ modulo $31$ divides the displayed defining resultant;
this is verified by \path{verification/verify_f4_mod31.sing}.

\subsection{Exact certificate and the known degree minimum}
\label{sec:certificate}

Write
\[
 P(T,W)=\sum_{j=0}^{23}p_j(T)W^j.
\]

\begin{proposition}[Exact function-field identity]\label{prop:identity}
In $\QQ(T)[V]/(\Phi)=\QQ(T)[V]/(\widehat F)$ one has
\[
 \sum_{j=0}^{23}p_j(T)J_0^jJ_1^{23-j}=0.
\]
Equivalently, $P(T,J_0/J_1)=0$ in $\QQ(X)$.
\end{proposition}

\begin{proof}
This is a polynomial identity modulo $\Phi=F$, equivalently modulo the monic
degree-$23$ relation $\widehat F$ over $\QQ(T)$.  The supplied Singular
certificate evaluates the left side by homogeneous Horner reduction and
obtains zero.  The SageMath certificate independently rebuilds
$\Phi,P,J_0,J_1$ from JSON coefficient tables, reads no Singular inputs, and
performs the same calculation in its own polynomial arithmetic.  A
self-contained Magma certificate repeats the quotient-ring identity from the
same canonical tables.  All computations are exact.
\end{proof}

\begin{proposition}[Irreducibility]\label{prop:irreducible}
The polynomial $P$ is primitive and irreducible in $\QQ[T,W]$.  In
particular it is irreducible in $\QQ(T)[W]$.
\end{proposition}

\begin{proof}
The coefficient content is $1$.  Its reduction modulo $31$ retains bidegree
$(23,4)$ and is irreducible in $\FF_{31}[T,W]$, as checked independently by
Singular, SageMath, and Magma.  Gauss's lemma proves the claim.
\end{proof}

\begin{proof}[Proof of Theorem~\ref{thm:main}]
By \cref{prop:identity}, $\QQ(T,W)$ is a subfield of $\QQ(T,V)$.  By
\cref{prop:irreducible},
\[
 [\QQ(T,W):\QQ(T)]=23.
\]
The source equation also gives $[\QQ(T,V):\QQ(T)]=23$, hence the inclusion
is an equality.  The regular $M_{23}$ monodromy follows from
\cite[Proposition~3.5 and Corollary~3.6]{HJLPPZ}.

Regard $P$ as a polynomial in $T$ over $\QQ(W)$.  Since $P$ is irreducible
and has positive $T$-degree, Gauss's lemma and $\deg_TP=4$ give
$[\QQ(T,W):\QQ(W)]=4$.  Thus $W$ has degree $4$ on $X$.  That this degree is
minimal over $\QQ$ is the nonsplit-quadric argument already given in
\cite[Remark~3.4]{HJLPPZ}: every map of degree at most $3$ on the geometric
genus-$4$ curve comes from a ruling of its canonical quadric, but neither
ruling is defined over $\QQ$.  Our calculation in
\cref{prop:canonical-ruling-field} identifies the quadratic ruling field
explicitly.
\end{proof}

\section{Arithmetic of the distinguished branch fibre}
\label{sec:branch-fibre-arithmetic}

The minimal equation and the original plane equation describe the same
cover, but their coordinates need not display a singular branch fibre in the
same way.  We first transport the fibre over $T=0$ exactly between the two
models.  We then determine its local monodromy, Fano-plane incidence, number
fields, and Galois-selected flags, before asking how much of this internal
arithmetic can explain the distinguished point of the seven-cover moduli
problem.

\subsection{Transport to the minimal-degree coordinate}
\label{sec:coordinate-transport}

The factors in the source coordinate are
\[
\begin{aligned}
H_7(V)={}&V^7-312V^6+43848V^5-3914896V^4
 +241123008V^3\\
&-9876755712V^2+251799776256V-2694891036672,\\
R_8(V)={}&V^8+156V^7-1290V^6-1217744V^5-40611765V^4\\
&+2044635024V^3+115083889782V^2
 +968032514052V-6058197515007.
\end{aligned}
\]
Both are irreducible over $\QQ$.  The corresponding primitive
factors in the minimal-degree coordinate are
\[
\begin{aligned}
\widetilde H_7(W)={}&52590311992108774196155W^7
-8380256819525975478304W^6\\
&+554571400016133612672W^5-19338639702208327424W^4\\
&+367883711159073536W^3-3404298821468160W^2\\
&+7203244769280W+62710185984,
\end{aligned}
\]
and
\[
\begin{aligned}
\widetilde R_8(W)={}&10551428180783458291047W^8
-2504505235463133310560W^7\\
&+258988485843116269632W^6-15238022132962377728W^5\\
&+557642314807148032W^4-12985270524477440W^3\\
&+187639106142208W^2-1535652528128W+5439225856.
\end{aligned}
\]

\begin{proof}[Proof of Theorem~\ref{thm:branch-fibre-transport}]
Reconstruct $F,P,J_0,J_1$ from their canonical integer coefficient tables.
Exact factorization over $\QQ$ gives the two displayed identities and proves
the irreducibility of all four factors.  Euclidean reduction gives
\[
 \gcd(J_1(0,V),H_7R_8)=1,
\]
so the specialized rational function is defined at every geometric point of
the fibre.

Let $v_7$ be the image of $V$ in $K_7=\QQ[V]/(H_7)$ and put
\[
 w_7=J_0(0,v_7)/J_1(0,v_7).
\]
Exact minimal-polynomial computation gives
\[
 \operatorname{minpoly}_{\QQ}(w_7)
 =\frac{\widetilde H_7}{52590311992108774196155}.
\]
In particular $w_7$ has degree $7$ and generates $K_7$.  Substitution
therefore defines the first displayed field isomorphism.  The identical
calculation in $K_8=\QQ[V]/(R_8)$ gives
\[
 \operatorname{minpoly}_{\QQ}\left(
 \frac{J_0(0,v_8)}{J_1(0,v_8)}\right)
 =\frac{\widetilde R_8}{10551428180783458291047},
\]
and proves the second isomorphism.  Applying all embeddings into
$\overline\QQ$ shows that the maps on roots are Galois-equivariant and
bijective.  Every step is reproduced by the certificate described in
\cref{sec:certificates}.
\end{proof}

\begin{remark}
The factor degrees provide a consistency check.  The theorem also requires
regularity of the rational coordinate throughout the fibre and generation of
the full residue fields.  The coprimality and minimal-polynomial checks verify
these two properties.
\end{remark}

\subsection{Local monodromy at the branch point
  \texorpdfstring{$T=0$}{T=0}}
\label{sec:setup}

\subsubsection{Group-theoretic notation}

Fix a copy of $\Mtt\subset\Sym(23)$ and elements
\[
 x\in2A,\qquad y\in23A,\qquad xy\in23A,
\]
so that $(x,y,(xy)^{-1})$ is a generating triple with class tuple
$(2A,23A,23B)$.  Set
\[
 K=O_2(\Cx),\qquad
 H=\Fix(x)\subset\{1,\dots,23\}.
\]
Since $x$ has cycle shape $1^72^8$, $|H|=7$; the fixed letters form a
heptad in the Witt design $S(4,7,23)$.  One has
\[
 \Cx\cong2^4{:}\PSL_3(2),\qquad |\Cx|=2688,
 \qquad |K|=16.
\]

For reproducibility, the certificate labeling uses
\[
\begin{aligned}
x={}&(3\,21)(4\,16)(9\,22)(10\,15)(11\,20)(12\,19)(13\,18)(14\,17),\\
y={}&(1\,2\,11\,10\,16\,9\,6\,3\,23\,19\,20\,14\,21\,17\,4\,8\,22\,5\,18\,15\,13\,7\,12),\\
r={}&(3\,18)(4\,16)(5\,6)(7\,23)(10\,14)(12\,19)(13\,21)(15\,17).
\end{aligned}
\]
The involution $r$ will compare the displayed triple with its complex
conjugate in \cref{thm:complex-conjugation}.  These representatives generate the same
Nielsen class as the tuple given in \cite[Section~3]{HJLPPZ}; all assertions below
are invariant under simultaneous conjugation.

For fixed $y$ let
\[
 X_y=\{x'\in2A:x'y\in23A\},\qquad |X_y|=161=7\cdot23,
\]
whose $\langle y\rangle$-translation orbits are the seven Nielsen classes.
We call the class containing $x$ the \emph{distinguished class}; it is
Nielsen class $6$ in the enumeration used below.  This is an internal label
for the class specified by the displayed permutation $x$.  The exact Hurwitz
algebra calculation in \cref{sec:hurwitz-scheme} later identifies
it intrinsically as the unique degree-one component over $K_0$.

Put $K_0=\QQ(\sqrt{-23})$.  Let $\mathcal H_{\mathrm{in}}$ denote the
reduced zero-dimensional inner Hurwitz scheme over $K_0$ for normalized
three-point covers with branch-cycle classes $(2A,23A,23B)$.  Here
\emph{inner} means that generating triples are taken modulo simultaneous
conjugation by $M_{23}$; we use the standard Hurwitz-space and Nielsen-class
formalism of \cite{FriedVolklein}.  Its seven geometric points correspond to the
seven $\langle y\rangle$-orbits above.  In
\cref{sec:hurwitz-scheme} we compute a finite \(K_0\)-scheme from exact
equations and then identify it with $\mathcal H_{\mathrm{in}}$ by certified
branch cycles.

\subsubsection{Branch points and the fibre over \texorpdfstring{$T=0$}{T=0}}

The cover is branched over $0$ and the two roots of $T^2+23$, with local
monodromy $(2A,23A,23B)$ and total ramification at the order-$23$ points
\cite[Section~3]{HJLPPZ}.  The valuations and ordinary nodes were computed
in \cref{lem:valuations,comp:nodes}.  The discriminant separates the branch
fibres from the projection singularities.

\begin{proposition}\label{thm:branch}
The integral model $F\in\ZZ[T,V]$ satisfies
\[
 \Res_V(F,F_V)=uT^8(T^2+23)^{92}G(T)^2,
 \qquad
 \disc_V(F)=-uT^8(T^2+23)^{88}G(T)^2,
 \qquad \deg G=84,
\]
for some $u\in\QQ^\times$, where $G$ is squarefree and divides
$\gcd(\Res_V(F,F_V),\Res_V(F,F_T))$, while $T$ does not.  Consequently:
\begin{enumerate}
\item the fibre over $T=0$ consists of smooth points of the plane model,
  and $F(0,V)=23^4H_7R_8^2$ gives cycle type $1^72^8$, the shape of $2A$;
\item the zeros of $G$ account for projection singularities rather than
  branch points of the normalization; their good reduction modulo $31$ is
  the $84$-node scheme of \cref{comp:nodes};
\item at $T^2+23$ the Newton polygon is a single segment of slope $4/23$,
  so the local extension is totally ramified of degree $23$.
\end{enumerate}
Riemann--Hurwitz gives
\[
 2g-2=-46+8+22+22=6,
 \qquad g=4.
\]
\end{proposition}

For comparison, direct factorization for the minimal-degree equation gives,
up to a
nonzero rational constant,
\[
 \disc_WP=T^8(T^2+23)^{22}H_{62}(T)^2,
 \qquad \deg H_{62}=62.
\]
The factor $T^8$ records genuine ramification in both models.  The residual
square factors and the exponent at $T^2+23$ depend on the projection.

\subsubsection{Why the cyclotomic obstruction vanishes}

\begin{proposition}\label{prop:qr}
$y^\lambda$ is conjugate to $y$ in $\Mtt$ if and only if $\lambda$ is a
quadratic residue modulo $23$; for the eleven nonsquares it lies in the
opposite class $23B$.
\end{proposition}

The two order-$23$ branch points are conjugate over $\QQ(\sqrt{-23})$, the
quadratic subfield of $\QQ(\zeta_{23})$ fixed by the squares.  Thus the
branch swap matches the nonsquare part of the cyclotomic character.  Fried's
branch-cycle argument sends $(C_1,C_2,C_3)$ under $\sigma$ to
$(C_1^\lambda,C_2^\lambda,C_3^\lambda)$, with
$\lambda=\chi(\sigma)$.  For nonsquare $\lambda$ the tuple becomes
$(2A,23B,23A)$ while the same $\sigma$ exchanges the two nonrational branch
points.  The classical obstruction therefore vanishes, but this permits
field of moduli $\QQ$ without forcing the fixed point of \cite{HJLPPZ}.

\subsection{The Fano plane on the seven unramified points}
\label{sec:fano-plane}

The factors $H_7$ and $R_8$ were displayed in
\cref{sec:coordinate-transport}.  Their
roots are respectively the seven unramified points and eight ramification
points above $T=0$.

\begin{theorem}\label{thm:fano-plane}
$\Gal(H_7)\cong\PSL_3(2)$, acting on the seven unramified sheets.  The $21$
involutions each fix exactly three sheets.  There are seven distinct fixed
triples, each arising from three involutions, and these triples form a
$2$-$(7,3,1)$ design.  Thus the sheets are the points of a Fano plane and
the fixed triples are its lines.
\end{theorem}

\smallskip
A \emph{flag} means an incident point--line pair in this plane.

\begin{proposition}\label{prop:linefield}
Let $\alpha_1,\ldots,\alpha_7$ be the roots of $H_7$.  The three-subset-sum
resolvent
\[
 \mathcal R_3(Y)=\prod_{1\leq i<j<k\leq7}
 \bigl(Y-(\alpha_i+\alpha_j+\alpha_k)\bigr)
\]
factors over $\QQ$ into irreducible factors of degrees $7$ and $28$.  The
degree-seven factor is
\[
\begin{aligned}
L_7(Y)={}&Y^7-936Y^6+379728Y^5-86588640Y^4
 +11922182400Y^3\\
&-982633068288Y^2+44572755885824Y-859311381083136.
\end{aligned}
\]
Its roots are the sums along the seven Fano lines.  The fields
$\QQ[V]/(H_7)$ and $\QQ[Y]/(L_7)$ are nonisomorphic, but they have the same
Galois closure and Dedekind zeta function.  Their common field discriminant
is
\[
 2^4 3^6 23^6=1726690609296.
\]
Their point and line stabilizers are nonconjugate subgroups of order $24$
with equal permutation characters, the classical Gassmann pair for
$\PSL_3(2)$; the resulting equality of Dedekind zeta functions is the
Gassmann--Perlis criterion \cite{Perlis}.
\end{proposition}

\begin{corollary}\label{cor:fano}
The roots of $\widetilde H_7$ carry the transported Fano-plane point
structure.  Its line field is defined by $L_7$, and the point and line fields
remain nonisomorphic arithmetically equivalent septics.
\end{corollary}

\begin{proof}
The first isomorphism of \cref{thm:branch-fibre-transport} identifies the
two seven-element
Galois sets.  Fixed triples and incidence are functorial under this
identification.
\end{proof}

\smallskip
\begin{proposition}\label{prop:centralizer-fano}
The points, lines, and flags admit the following description through
$\Cx$.  All $99$ involutions of $\Cx$ belong to class $2A$,
and:
\begin{enumerate}
\item the fifteen nonidentity elements of $K$ fix $H$ pointwise.  The
  central element $x$ is one conjugacy class and the other fourteen form a
  second.  Writing $U=K/\langle x\rangle\cong\FF_2^3$, the nonzero elements
  of $U$ identify with the seven lines; the plane is $\PP(U^*)$ on points
  and $\PP(U)$ on lines;
\item as an $\FF_2[\Cx]$-module, $K$ is indecomposable, with unique minimal
  nonzero submodule $\langle x\rangle$.  Each nonzero element of $U$ has two
  preimages in $K$, exchanged by multiplication by $x$;
\item the $84$ involutions of $\Cx\smallsetminus K$ form one class and act
  on $H$ as transvections.  The map to center and fixed line is a surjection
  onto the $21$ flags, with four involutions above each flag.  For such $t$,
  \[
   C_K(t)=\langle x\rangle\cup\{\text{the two lifts of the axis}\},
   \qquad \Fix(t)\cap H=\operatorname{axis}(t).
  \]
\end{enumerate}
\end{proposition}

There is no $\Cx$-invariant complement to $\langle x\rangle$ in $K$.
Moreover the outer point--line duality of $\PSL_3(2)$ is not induced by an
automorphism of $\Cx$ or by a permutation of the $23$ sheets normalizing
$\Cx$:
\[
 \Aut(\Cx)=\Inn(\Cx)\quad(\text{order }1344),
 \qquad N_{\Sym(23)}(\Cx)=\Cx.
\]
Thus no such symmetry exchanges the point and line Galois sets.

\subsection{Number fields obtained from the fibre}
\label{sec:field}

Put
\[
 E=\Spl_{\QQ}(H_7),\qquad L_0=\Spl_{\QQ}(H_7R_8).
\]
We construct a third field $M$ and obtain
\[
 \QQ\subset E\subset L_0\subset M,
 \qquad [E:\QQ]=168,\quad[L_0:E]=8,\quad[M:L_0]=2.
\]
The successive groups over $\QQ$ are $\PSL_3(2)$,
$2^3{:}\PSL_3(2)$, and $2^4{:}\PSL_3(2)$.  Thus the three stages remember,
respectively, the seven unramified points, the eight ramification points,
and the choice between the two local branches at each ramification point.

\subsubsection{The splitting field of \texorpdfstring{$H_7R_8$}{H7R8}}

The local decomposition group centralizes $x$.  Since $x$ exchanges the two
local branches at a ramification point but fixes its root of $R_8$, the
action on the eight roots factors through $\Cx/\langle x\rangle$.  The
translation subgroup is $U=K/\langle x\rangle\cong\FF_2^3$.
After choosing an affine origin, the difference of two roots is a nonzero
vector of $U$, hence a Fano line by \cref{prop:centralizer-fano}; each of the
seven nonzero differences occurs for four of the $28$ unordered pairs.

\begin{theorem}\label{thm:affine}
The splitting field of $H_7R_8$ has Galois group
$2^3{:}\PSL_3(2)$ of order $1344$, acting affinely on the eight roots of
$R_8$.  The restriction $\Gal(L_0/\QQ)\to\Gal(E/\QQ)$ has kernel
$K/\langle x\rangle\cong(\ZZ/2)^3$.
\end{theorem}

\begin{proof}
Decomposition theory at the tame rational place, together with equality of
arithmetic and geometric monodromy \cite{HJLPPZ}, embeds
$\Gal(L_0/\QQ)$ in $\Cx/\langle x\rangle$.  Restriction to $E$ is
surjective, and its kernel is a submodule of the irreducible
three-dimensional module $U$, hence either trivial or all of $U$.  At
$p=599$, $H_7$ and $L_7$ split completely while $R_8$ factors into four
quadratics.  Frobenius is therefore trivial on $E$ but a nonzero translation
on the roots of $R_8$, proving that the kernel is $U$.
\end{proof}

Independently, a direct number-field computation gives
$[\Spl_{\QQ}(R_8):\QQ]=1344$.

\subsubsection{Affine frames and specialization}

Let $A$ be the eight roots of $R_8$.  An ordered affine frame is a tuple
$(a_0,a_1,a_2,a_3)$ whose three differences $a_i-a_0$ form a basis of $U$.
Let
\[
 \mathcal F(A)=\{\text{ordered affine frames of }A\}/U.
\]

\begin{proposition}\label{prop:affineframes}
$\mathcal F(A)$ has $168$ elements.  The action of
$\Cx/\langle x\rangle$ descends to a free transitive action of
$\Cx/K\cong\PSL_3(2)$.  Via \cref{prop:centralizer-fano}, a translation
class of frames is equivalently an ordered basis of the Fano line module.
\end{proposition}

\begin{proof}
There are $8\cdot7\cdot6\cdot4=1344$ affine frames.  Translation acts
freely on the eight origins, leaving $168$.  The quotient
$\PSL_3(2)=\operatorname{GL}_3(2)$ acts freely and transitively on ordered
bases of $U$.
\end{proof}

Let $S$ be the marked geometric generic fibre, based at the tangential
basepoint $T\to0$.  For $\sigma\in\GQ$, let $d(\sigma)$ be its descent
permutation.  Since the inertia $x$ has order two, the branch-cycle lemma
places $d(\sigma)$ in $\Cx$.  Write
$\rho_8:\GQ\to\Sym(A)$ for the Galois action on the roots of $R_8$.

\begin{proposition}[Tame specialization at $T=0$]
\label{prop:specialization}
For every $\sigma\in\GQ$,
\[
 \rho_8(\sigma)=\operatorname{pair}(d(\sigma)),
\]
where $\operatorname{pair}$ is the induced action on the eight two-sheet
inertia orbits.  Consequently the induced action on $\mathcal F(A)$ is the
one obtained from the generic-fibre descent permutation.
\end{proposition}

\begin{proof}
Let $R=\QQ[T]_{(T)}^h$ and normalize $\operatorname{Spec}R$ in the cover.
At a geometric point of ramification index $e$, tameness and smoothness give
the completed local form $T=u^e$ after extracting an $e$th root of a unit.
The $e$ sheets form one inertia orbit and specialize to one point.  Hence
specialization gives a canonical Galois-equivariant bijection
\[
 S/\langle x\rangle\xrightarrow{\sim}
 \{\text{geometric points above }T=0\}.
\]
Restricting to the eight two-element orbits proves the formula.  This is the
finite-set case of tame specialization
\cite[Expos\'e XIII, Section~2.10]{SGA1}.
\end{proof}

\subsubsection{Separating the two branches}

The group $\Cx$ is perfect, its center is $\langle x\rangle$, and the
central extension
\[
 1\longrightarrow\langle x\rangle\longrightarrow\Cx
 \longrightarrow\Cx/\langle x\rangle\longrightarrow1
\]
does not split.  It acts faithfully and transitively on the sixteen local
branches above the eight roots of $R_8$, with branch stabilizer isomorphic
to $\PSL_3(2)$.

Near a root $v$ of $R_8$, the two branches have expansions
\[
 V=v\pm\sqrt{e(v)T}+O(T),\qquad
 e(v)=-\frac{2F_T(0,v)}{F_{VV}(0,v)}.
\]
Since $F_{VV}(0,v)=2\cdot23^4H_7(v)R_8'(v)^2$, the square class of $e(v)$ is
represented by
\[
 c(V)\equiv-F_T(0,V)H_7(V)\pmod{R_8}.
\]

\begin{theorem}\label{thm:field}
The eight values $c(v)$ represent one nontrivial square class in
$L_0^\times/L_0^{\times2}$.  If $M=L_0(\sqrt{c(v)})$ for any root $v$, then
\[
 \Gal(M/\QQ)\cong\Cx=2^4{:}\PSL_3(2),
 \qquad [M:L_0]=2,
\]
and the nontrivial element of $\Gal(M/L_0)$ is $x$.
\end{theorem}

\begin{proof}
Tame local theory identifies the splitting field over $\QQ((T))$ as
$\mathcal L=L_0((\sqrt{cT}))$, with Galois group $\Cx$.  Since $\Cx$ is
perfect, $\sqrt T\notin\mathcal L$ and $c$ is not a square in $L_0$.  After
adjoining $\sqrt T$ the Galois group is $\Cx\times\ZZ/2$.  The diagonal
involution $(x,-1)$ fixes $\sqrt c=\sqrt{cT}/\sqrt T$, so taking constant
fields gives
\[
 \Gal(M/\QQ)\cong
 (\Cx\times\ZZ/2)/\langle(x,-1)\rangle\cong\Cx.
\]
Exact arithmetic checks additionally show that $X^2-dc(V)$ is irreducible
over $\QQ[V]/(R_8)$ for
\[
 d\in\{\pm1,\pm2,\pm3,\pm6,\pm23,\pm46,\pm69,\pm138\}.
\]
At each of fourteen primes below $3\cdot10^5$ splitting completely in
$L_0$, the eight values have one quadratic character; at six, beginning
with $13049$, it is $-1$, independently certifying nonsquareness.
\end{proof}

The polynomial $R_8$ has four real roots.  Exactly two have $c(v)>0$, so
both local branches above them are real; the other two have $c(v)<0$, so
their branches are exchanged by complex conjugation.  This agrees with the
permutation $r$, which fixes the branches above two real ramification points
and exchanges the branches above the other two.

\begin{corollary}\label{cor:affine}
Replacing $(H_7,R_8)$ by $(\widetilde H_7,\widetilde R_8)$ leaves the fields
$E,L_0,M$, the affine torsor, and the centralizer action unchanged under the
transport induced by $W=J_0/J_1$.
\end{corollary}

\begin{proof}
Theorem~\ref{thm:branch-fibre-transport} identifies both residue Galois
sets.  The first two
splitting fields and their functorial actions agree.  The final quadratic
extension belongs to the normalized cover and separates its local branches,
independently of the coordinate labeling the residue points.
\end{proof}

\subsubsection{Ramification}

\begin{theorem}\label{thm:ramification}
Let $E_8=\QQ[V]/(R_8)$ and $E'_8=E_8(\sqrt{c(v)})$.  The fields
$\QQ[V]/(H_7)$, $\QQ[Y]/(L_7)$, and $E_8$ have discriminants
$2^43^623^6$, $2^43^623^6$, and $2^{16}3^623^6$.  The relative
discriminant of $E'_8/E_8$ has norm $2^83^223^2$.  Thus $M$ is unramified
outside $\{2,3,23\}$ and $M/L_0$ ramifies above all three primes.  More
precisely:
\begin{enumerate}
\item at $23$, the septic fields are totally ramified with $e=7$.  Residue
  inertia is a Singer cycle, transitive on points and lines and fixing one
  affine origin.  The residue decomposition group is $7{:}3$; in $M$,
  inertia is $C_{14}$ and decomposition is $C_2\times(7{:}3)$;
\item at $3$, ramification is wild with $e=3$ and lies in the
  $\PSL_3(2)$ quotient;
\item at $2$, the septics are tame with $e=3$, while the octic is wild with
  $e=4$.  The relevant order-$12$ subgroup of $\Cx$ is $A_4$, with Klein
  four wild part inside $K$ and tame quotient $C_3$.
\end{enumerate}
\end{theorem}

\begin{proof}
The absolute and relative discriminants are computed from exact maximal
orders.  At $23$, tame local theory gives
$\phi\tau\phi^{-1}=\tau^{23}=\tau^2$; since $2$ has order three modulo $7$,
the residue decomposition group is $7{:}3$.  The quadratic step adjoins $x$
to inertia.  The assertions at $2$ and $3$ follow from the exact
factorization and ramification data for $H_7,L_7,R_8$ and the subgroup
census in the certificates.
\end{proof}

\subsection{Flags selected by complex conjugation and Frobenius}
\label{sec:flags}

\begin{theorem}\label{thm:flags}
Every involution of $\PSL_3(2)$ acts as a transvection.  It fixes one line
pointwise, its axis, and has a distinguished center on that line.  The map
from involutions to center--axis pairs is a bijection onto the $21$ flags.

On points the cycle type is $1^32^2$.  Besides the axis, exactly two lines
are invariant, each consisting of the center and one transposed pair.  The
group is transitive on flags; both a flag stabilizer and an involution
centralizer have order $8$.

The field $E$ has no real embeddings and $84$ conjugate pairs of complex
embeddings.  Each complex place determines the flag of its conjugation
involution, giving a four-to-one map from the $84$ places to the $21$ flags.
The same construction applies at an unramified finite place whose Frobenius
is an involution.
\end{theorem}

\begin{corollary}\label{cor:archimedean}
For complex conjugation, the axis is the set of three real roots of $H_7$.
Their sum is a root of $L_7$.  The center is the unique real root lying on
the other two invariant lines; numerically it is the largest real root,
approximately $102.2684$.
\end{corollary}

If $p$ is unramified and $H_7$ has factorization type $1^32^2$ modulo $p$,
then Frobenius is an involution.  Such primes have Chebotarev density $1/8$.

\begin{theorem}\label{thm:transvection-primes}
Let $N_{28}$ be the degree-$28$ nonline factor of $\mathcal R_3$, and
exclude primes dividing $\Res(L_7,N_{28})$.  At every remaining unramified
prime where $H_7$ has type $1^32^2$, the three roots in $\FF_p$ sum to a
root of $L_7$ and form the axis.  The other four roots form two conjugate
pairs.  Counting, for each fixed root, how many pairs complete a Fano line
gives $(0,0,2)$; the root with count $2$ is the center.  Exact computation
checks all $9{,}889$ such primes below $10^6$.  Below that bound the only
excluded collision is $p=22159$.
\end{theorem}

The real axis is defined over $\mathbb R$ but not over $\QQ$, since
$\Gal(E/\QQ)$ is transitive on the seven lines.  At $23$ one sees the
complementary Singer-cycle picture: after cyclically labeling the seven
points, one Fano line is a difference set and the other six are its
translates.

\subsection{Complex conjugation and the geometric reflection obstruction}
\label{sec:complex-conjugation}

The branch divisor is preserved by the rational base involution
\[
 \tau:\PP^1_T\longrightarrow\PP^1_T,\qquad T\longmapsto-T,
\]
which fixes the $2A$ branch point and exchanges the two order-$23$ branch
points.  It does not, however, lift to any of the seven covers.  The two
endpoint inertia generators lie in the distinct classes $23A$ and $23B$.
A holomorphic branch swap preserves the orientation of a small loop and
therefore preserves its inertia conjugacy class; inversion occurs only for
an orientation-reversing comparison such as complex conjugation.

\begin{proposition}[Geometric reflection obstruction]
\label{prop:arithmetic-reflection}
The holomorphic involution $\tau:T\mapsto-T$ has no lift to a cover in the
Nielsen class $(2A,23A,23B)$.  In contrast, the orientation-reversing
comparison equations
\[
 r\in C_{M_{23}}(x),\qquad
 y^r=z^{-1},\qquad z^r=y^{-1}
\]
have solution counts $(0,0,1,0,0,1,1)$ on the seven normalized triples.
Thus IDs $3,6,7$ are the classes fixed by the chosen complex-conjugation
action.
\end{proposition}

\begin{proof}
For the natural degree-$23$ action, exact computation gives
\[
 C_{S_{23}}(M_{23})=1,
 \qquad N_{S_{23}}(M_{23})=M_{23}.
\]
For every normalized triple, $y$ is not conjugate to $z$ in $M_{23}$,
whereas it is conjugate to $z^{-1}$.  Any sheet permutation induced by a
holomorphic lift would normalize $M_{23}$ and carry a positive local
inertia generator to a conjugate positive generator at the other endpoint.
The normalizer identity would place it in $M_{23}$, contradicting the first
class calculation.  Complex conjugation reverses the local loop, so its
finite comparison uses $z^{-1}$ instead.  Exhausting the displayed
orientation-reversing equations gives the asserted counts.
\end{proof}

Let
\[
 Y_x=\{y''\in23A:xy''\in23A\}.
\]
Pairs $(x,y'')$ represent the seven Nielsen classes, and $\Cx$ acts on
$Y_x$ by conjugation.

\begin{proposition}\label{prop:torsor}
The group $\Mtt$ acts freely on the admissible pairs
$(x',y')$ with $x'\in2A$, $y'\in23A$, and $x'y'\in23A$.  Moreover
\[
 |Y_x|=18816=7\cdot2688,
\]
and $Y_x$ is the disjoint union of seven free $\Cx$-orbits, one per Nielsen
class.  Equivalently,
\[
 |Y_x|=|X_y|\frac{|23A|}{|2A|}
 =161\cdot\frac{443520}{3795}=18816.
\]
\end{proposition}

In these coordinates complex conjugation is
\[
 \mu:Y_x\longrightarrow Y_x,\qquad \mu(y'')=xy''.
\]

\begin{theorem}\label{thm:complex-conjugation}
The equation
\[
 y''{}^w=xy'',\qquad w\in\Cx,
\]
has a solution exactly when the corresponding Nielsen class is fixed by
complex conjugation.  Exactly three classes are fixed.  For each, the
comparison element $w$ is unique, lies in $\Cx\smallsetminus K$, and acts as
a transvection on the roots of $H_7$.  For the distinguished class, the unique
solution is $w=r$; it satisfies $x^r=x$ and $y^r=xy$.  Its center and axis
give the real flag of \cref{thm:flags}.
\end{theorem}

At the chosen real place, complex conjugation exchanges the same two branch
points and induces this orientation-reversing permutation on the Nielsen set.
The element $r$ acts on the roots of $H_7$ with type $1^32^2$, fixes three
Fano lines, and fixes four roots of $R_8$.  Correspondingly, $H_7,L_7,R_8$
have $3,3,4$ real roots.  The three fixed classes are the fixed set of this
one complex-conjugation element.  A $\QQ$-defined cubic subscheme would
require stability under the full absolute Galois group.  The calculation
neither establishes such stability nor explains why the distinguished class
is globally fixed.

To relate the fibre arithmetic to the fixed-class problem, we now record the
information supplied by the branch-fibre structures and the additional
comparison that remains missing.

\subsection{From the branch fibre to the fixed Nielsen class}
\label{sec:limits}

The construction of \cite{HJLPPZ} starts from a seven-element Nielsen class.
Its Galois-fixed point produces the rational cover studied here.  The Fano
plane is attached instead to the seven unramified points in the chosen $2A$
branch fibre of that cover.  Thus the results above start from the known
rational cover and describe arithmetic internal to it; its descent is an
input from \cite{HJLPPZ}.

The roles of the structures above may be summarized as follows.
\begin{enumerate}
\item The minimal function $W$ transports the residue fields between the two
  plane models, while the Fano incidence is intrinsic to the chosen fibre.
\item The Fano plane is intrinsic to the seven fixed sheets of the chosen
  $2A$ fibre, but no individual Fano line is Galois invariant over $\QQ$.
\item The affine-frame torsor is compatible with tame specialization from
  the geometric generic fibre to the fibre over $T=0$.  Distinguishing the
  rational Nielsen class from the other six requires additional descent data.
\end{enumerate}

\smallskip
The coordinate comparison preserves all Fano and affine residue data, so a
future explanation based on the branch fibre may be formulated entirely in
the minimal-degree model.  The exact Hurwitz algebra computed later identifies
the degree-one component directly.  Constructing that component conceptually
from the branch-fibre data remains open.

\smallskip
There is also an exact, entirely group-theoretic limit on what unpointed
characteristic-$23$ inertia can determine.  Let $I\leq\Mtt$ have order $23$
and put $D=N_{\Mtt}(I)$.  Direct computation in the supplied permutation model
gives $D\cong23{:}11$.  The conjugates of the unpointed inertia subgroup $I$ are
therefore indexed by the $40320$ points of $\Mtt/D$.  Once a chosen local
object is reduced to its unpointed order-$23$ inertia subgroup, its possible
placements are exactly the points of this homogeneous space.

\begin{proposition}[Unpointed $23$-inertia obstruction]
\label{prop:gauss-obstruction}
For each of the three complex-conjugation-fixed Nielsen classes, let $r$ be
the orientation-reversing comparison of
\cref{prop:arithmetic-reflection} and let
\[
 A=C_{\Mtt}(x)\cap C_{\Mtt}(r).
\]
Then $|A|=32$, the three ordered pairs $(x,r)$ are simultaneously
conjugate in $\Mtt$, and $A$ acts freely on the $40320$ conjugates of $I$,
indexed by $\Mtt/D$.  The involution $r$ divides them into $20160$ unordered pairs,
on which $A/\langle r\rangle$ acts freely.  Consequently there are $1260$
orbits of reflection-pairs under $A/\langle r\rangle$ for each fixed
Nielsen class.
In particular, unpointed order-$23$ inertia together with $(x,r)$ does not
canonically select one pair or distinguish the distinguished class from the
 other two complex-conjugation-fixed classes.
\end{proposition}

\begin{proof}
Exact computation gives $|A|=32$ and one simultaneous-conjugacy orbit for
the three anchors.  Since $|D|=253$ is odd, the stabilizer in the
$2$-group $A$ of any coset in $\Mtt/D$ is an intersection with a conjugate
of $D$ and is therefore trivial.  Thus $A$ acts freely.  The element $r$
is central in $A$ and has no fixed coset.  An element of $A$ stabilizing
one of its unordered two-element orbits either fixes a coset, hence is the
identity, or sends it to its $r$-translate, hence is $r$.  The induced
$A/\langle r\rangle$-action is free, and
\[
 \frac{20160}{|A|/2}=\frac{20160}{16}=1260.
\]
\end{proof}

Thus a $23$-adic argument must augment the two inertia subgroups with
local--global comparison data.  More precisely, the characteristic-$23$
Kummer coordinate has two $C_{11}$-orbits on its
nonzero parameters: the squares in $\FF_{23}^{\times}$ and their nonsquare
coset.  On the Fano side, if $F_{\mathrm{fl}}=E^{D_8}$ is a flag field,
the two Klein four subgroups $V_4\subset D_8$ give two quadratic extensions
$E^{V_4}/F_{\mathrm{fl}}$.  A proposed explanation along these lines would
need a Galois-equivariant comparison between these two explicitly defined
two-element structures.

\smallskip
There is another natural quadratic datum on the curve, but it does not
supply this comparison.  Let $(A_0,A_1,A_2,A_3)$ be the RREF-normalized
basis of the canonical adjoint space reconstructed from the same $88$
coefficient positions used above.  Order quadratic monomials as
\[
 X_0^2,X_0X_1,X_0X_2,X_0X_3,X_1^2,X_1X_2,X_1X_3,
 X_2^2,X_2X_3,X_3^2.
\]

The nonsquare discriminant of this canonical quadric is the obstruction used
in \cite[Remark~3.4]{HJLPPZ} to prove that degree $4$ is minimal over $\QQ$.
The next proposition refines that argument by identifying the ruling field
exactly; that field is also relevant to the local--global comparison pursued
here.

\begin{proposition}[Canonical ruling field]
\label{prop:canonical-ruling-field}
The unique quadric containing the canonical image, normalized to have
coefficient $1$ at $X_3^2$, has coefficient vector
\[
\frac1{297}\bigl(
\begin{gathered}
442260599149123,-18897532752913,-34200754149,794833324,\\
168493762065,745707549,-14145827,660893,-31345,297
\end{gathered}
\bigr).
\]
If $B$ is twice its symmetric Gram matrix, then
\[
 \det B=\frac{644454138716416151027888}{970299}
       =4873\left(\frac{38141181796}{3267}\right)^2,
 \qquad 4873=11\cdot443.
\]
Consequently the two rulings are governed by
\[
 \QQ(\sqrt{4873}),
\]
which is different from the branch-orientation field
$\QQ(\sqrt{-23})$.
\end{proposition}

\begin{proof}
The four canonical RREF vectors are reconstructed over $\QQ$ from $39$
good primes with a $604$-bit CRT modulus; the ten quadric coefficients use
a $141$-bit modulus.  Both reconstructions agree coefficient by coefficient
at the independent holdout prime $137$, and the canonical basis reduces to
the directly recomputed basis modulo $31$.  Exact reduction in
$\QQ(T)[V]/(\widehat F)$ annihilates the displayed quadratic combination
of the ten products $A_iA_j$; their coefficient matrix has rank $9$, proving
uniqueness.  Exact determinant evaluation gives the displayed square class.
The components of the Fano scheme of lines on a smooth quadric surface are
defined over its discriminant extension, so the ruling field is
$\QQ(\sqrt{4873})$.  The SageMath and Magma certificates independently
repeat the quotient-ring, rank, and determinant calculations.
\end{proof}

The quadratic extension governing the two rulings cannot supply this
comparison.  Indeed, retain the notation
$F_{\mathrm{fl}}=E^{D_8}$ and let $E^{V_4}/F_{\mathrm{fl}}$ be either of
the two quadratic flag extensions just defined.  If its square class were
the product of the square classes obtained from the branch-orientation
field $\QQ(\sqrt{-23})$ and the ruling field $\QQ(\sqrt{4873})$, then
\[
 E^{V_4}=F_{\mathrm{fl}}\QQ(\sqrt{-23\cdot4873}).
\]
The right side would imply
$\QQ(\sqrt{-23\cdot4873})\subset E$.  This is impossible because
$\Gal(E/\QQ)\cong\PSL_3(2)$ has no quotient of order two.  Thus the ruling
field does not provide the required local--global comparison.

The reflection, inertia, and ruling-field calculations determine the reach
and the limits of the arithmetic internal to the distinguished cover.  They
produce candidate comparison data and precise obstructions, but they do not
by themselves construct the degree-one component of the seven-cover moduli
problem.  We therefore leave the single branch fibre and compute the full
finite Hurwitz scheme and its Galois action directly.

\section{The finite Hurwitz scheme and the six conjugate maps}
\label{sec:hurwitz-scheme}

\subsection{Coordinate algebra and canonical models}

We now return from the $K_0$-rational point corresponding to the distinguished
class to the whole inner Nielsen class.  The normalized canonical quadric has the
normalization-invariant coefficient
\[
 J=\frac{q_{13}}{A^2},
\]
where $A$ is the coefficient used to impose the fixed normalization in the
accompanying data.  Computing $J$ at the seven analytic classes and
reconstructing over $K_0$ gives a finite étale $K_0$-algebra
\[
 A_{\mathrm H}=K_0\times L,
 \qquad [L:K_0]=6.
\]
Let $\mathcal H_J=\operatorname{Spec}A_{\mathrm H}$.  Its degree-one
component carries the reconstructed map belonging to the distinguished class,
and the six geometric points of its connected degree-six component carry the
other six
reconstructed maps.  The branch-cycle computation below identifies
$\mathcal H_J$ with the inner Hurwitz scheme $\mathcal H_{\mathrm{in}}$.

\begin{computation}[Coordinate algebra and canonical curves]
\label{comp:hurwitz-coordinate-algebra}
The linear factor of the polynomial of $J$ is
\[
 J-\frac{148227-34830\sqrt{-23}}{142129},
\]
and its complementary factor is irreducible of degree $6$ over $K_0$.
An absolute defining polynomial for $L$ is
\[
\begin{split}
 h(Z)={}&Z^{12}-6Z^{11}+20Z^{10}-32Z^9+44Z^8-22Z^7+6Z^6\\
       &{}-22Z^5+44Z^4-32Z^3+20Z^2-6Z+1.
\end{split}
\]
Over $A_{\mathrm H}$, all ten coefficients of the canonical quadric, all
twenty coefficients of a Petri cubic in the fixed coefficient
normalization, and the two marked
order-$23$ points reconstruct exactly.  On each component the resulting
quadric--cubic complete intersection is smooth.
\end{computation}

The reciprocal form of $h$ makes the Galois closure unexpectedly small to
describe.  Put
\[
 g(Y)=Y^6-6Y^5+14Y^4-2Y^3-27Y^2+44Y-44.
\]

\begin{proposition}[Galois closure of the sextic component]
\label{prop:hurwitz-galois-closure}
Let $E=\QQ[Y]/(g)$ and let $\widetilde E$ be its Galois closure.  Then
\[
 L=EK_0,\qquad
 \Gal(\widetilde E/\QQ)\cong S_6,qquad
 \Gal(\widetilde E K_0/K_0)\cong S_6,
\]
and
\[
 \Gal(\widetilde E K_0/\QQ)\cong S_6\times C_2.
\]
The relative Galois group acts as the full symmetric group on the six
remaining maps.  In the embedding order of
\cref{comp:hurwitz-branch-cycles}, its letters have Nielsen labels
\[
 (7,4,1,5,3,2).
\]
Moreover
\[
 \disc(E)=2^4\cdot11\cdot23^4,qquad
 \mathfrak d_{L/K_0}=(2^4\cdot11\cdot23).
\]
\end{proposition}

\begin{proof}
Let $a$ be the image of $Z$ in $\QQ[Z]/(h)$.
The reciprocal-polynomial identity is
\[
 h(Z)=Z^6g(Z+Z^{-1}).
\]
The exact relative-to-absolute field map sends $\sqrt{-23}$ into $\QQ(a)$,
and the element $a+a^{-1}+\sqrt{-23}$ has degree $12$ over $\QQ$.
Consequently $L=EK_0$.  Maximal-order calculations give the displayed
discriminants.

The rational primes $3,139,2671$ split in $K_0$, do not divide $\disc(g)$,
and give respective factor-degree patterns
\[
 (6),\qquad(1,5),\qquad(1,1,1,1,2)
\]
for $g$ modulo those primes.  Thus the relative Galois group contains a
$6$-cycle, a $5$-cycle, and a transposition.  The first factorization also
proves irreducibility.  Of the sixteen transitive subgroups of $S_6$, only
$6T16=S_6$ contains all three cycle types.

Finally, the unique quadratic subfield of an $S_6$ splitting field is its
discriminant field.  Since
\[
 \disc(g)=2^{22}\cdot11\cdot23^4,
\]
that field is $\QQ(\sqrt{11})$, rather than $K_0=\QQ(\sqrt{-23})$.
Therefore $\widetilde E$ and $K_0$ are linearly disjoint, giving the two
asserted Galois groups.  The Nielsen labeling is the exact embedding
crosswalk from the branch-cycle continuation.
\end{proof}

\subsection{Local structure at \texorpdfstring{$23$}{23}}

The prime above $23$ gives a particularly small local realization of this
full symmetric action.  Let $\mathfrak q=(\sqrt{-23})$ be the unique prime of
$K_0$ above $23$.

\begin{proposition}[Local Hurwitz algebra at $23$]
\label{prop:hurwitz-local-23}
The normalization of $\operatorname{Spec}\mathcal O_{K_0,\mathfrak q}$ in
$A_{\mathrm H}=K_0\times L$ has three closed points.  Their ramification
indices and residue degrees over $K_{0,\mathfrak q}$ are
\[
 (e,f)=(1,1),\qquad(1,2),\qquad(2,2),
\]
where the first point lies on the degree-one component and the other two lie on
the sextic component.  In the $S_6$-closure, a decomposition group at
$\mathfrak q$ is
\[
 D_{\mathfrak q}\cong V_4.
\]
Its natural action on the six remaining points has orbits of sizes $2$
and $4$; inertia has cycle type $1^2 2^2$, and either lift of residue
Frobenius has cycle type $2^3$.
\end{proposition}

\begin{proof}
Modulo $23$ the trace sextic satisfies
\[
 g(Y)=(Y-5)^2(Y+1)^4.
\]
Over $\QQ_{23}$ its two factors have degrees $2$ and $4$.  Their
discriminant valuations are $1$ and $3$, with residual units $9$ and $7$;
their discriminant square classes are therefore $23$ and $-23$.
The quartic factor is totally tamely ramified.  Its local Galois closure is
dihedral of order $8$: tame inertia is cyclic of order $4$, while Frobenius
acts by inversion because $23\equiv-1\pmod4$.  Its discriminant field is
$\QQ_{23}(\sqrt{-23})$.  Restriction to $K_{0,\mathfrak q}$ gives the even
Klein four subgroup.  The quadratic factor field is
$\QQ_{23}(\sqrt{23})$; its compositum with $K_{0,\mathfrak q}$ is the same
unramified quadratic field already contained in the dihedral closure.
This gives the asserted group action and prime data.

The first residual equations make the common residue field explicit.  If
$s=\sqrt{-23}$, then
\[
 U=\frac{Y-5}{s},\quad \overline U^{\,2}=15,
 \qquad
 V=\frac{(Y+1)^2}{s},\quad \overline V^{\,2}=5.
\]
Both equations define $\FF_{23^2}$, since $15/5=3=7^2$ in $\FF_{23}$.
\end{proof}

A complex Nielsen labeling of the local $2+4$ partition is not canonical:
it requires choosing a prime of the $S_6$-closure above $\mathfrak q$, and
changing that prime conjugates $D_{\mathfrak q}$ inside $S_6$.

We use the standard ADE nomenclature and the Milnor, Tjurina, and delta
invariants as in \cite{GLS}.

\begin{computation}[Pointed direct reductions at $23$]
\label{comp:hurwitz-pointed-23}
The exact canonical equations, marked points, and degree-$23$ map sections
are integral at all three local Hurwitz points.  The direct special fibres
are geometrically integral quadric--cubic complete intersections, their
canonical quadrics remain smooth, and the marked zero $b$ and pole $c$
remain distinct smooth points.  Their singularities are
\[
\begin{array}{c|c|c}
\text{local point}&\text{residue field}&\text{singularities}\\
\hline
\text{degree-one point (class 6)}&\FF_{23}&E_8\\
\text{unramified degree }2&\FF_{23^2}&E_8\\
\text{ramified degree }4&\FF_{23^2}&A_2+A_6.
\end{array}
\]
The Milnor and Tjurina numbers are respectively $8$, and $2+6$.  In each
row the total delta invariant is $4$, so the geometric normalization is
$\PP^1$.

Let $N,D$ be the reduced quintic map sections.  Their common base scheme
has length $7$, avoids the singular points, and
\[
 \ord_b(N)=\ord_c(D)=23,
 \qquad \ord_b(D)=\ord_c(N)=0.
\]
Consequently, on the normalization,
\[
 \operatorname{div}(N/D)=23b-23c.
\]
After choosing a coordinate $t$ with $b=0$ and $c=\infty$ and rescaling it,
the pointed direct reduction is
\[
 \overline\beta=t^{23}
\]
at all three local points.
\end{computation}

\begin{proof}
Exact reduction in the two residue fields gives the displayed singular
loci and no singularity at infinity.  Eliminating the smooth-quadric
coordinate gives an irreducible plane sextic.  The exact local plane
equations have
triple tangent, Milnor--Tjurina pair $(8,8)$, and fifth-order tangent
restriction in the first two cases, identifying $E_8$ in characteristic
$23$.  At the third point the two completed local equations have
Milnor--Tjurina pairs $(2,2)$ and $(6,6)$ and critical-value orders $3$ and
$7$, identifying $A_2$ and $A_6$.  Their delta invariants exhaust the
arithmetic genus $4$.  The common base ideal has constant Hilbert polynomial
$7$, and exact local series at $b,c$ give the four displayed orders.  Since
the residue fields are perfect, the remaining scalar multiplying $t^{23}$
can be absorbed into $t$.
\end{proof}

This calculation does not yet give the desired stable pointing.  Direct
normalization retains only the purely inseparable Frobenius map; the genus
and the separable Mathieu tail are concentrated in the ADE singularities.
One must resolve their mixed-characteristic deformations while retaining
the gluing to a common pointed $M_{23}$-frame.

It does, however, give the binary augmentation a concrete local arithmetic
interpretation.  If $P$ is one of the three closed points in
\cref{prop:hurwitz-local-23} and $f(P)$ is its residue degree, then under the
already certified Hurwitz-algebra identification,
\[
 \epsilon(\Theta)(P)=f(P)-1\pmod2.
\]
Indeed, both sides have values $(0,1,1)$ on the residue-degree pattern
$(1,2,2)$.  This equality uses the identified local Hurwitz algebra; it is
not yet an intrinsic construction of $\epsilon(\Theta)$ from stable maps.

\subsection{Reconstruction and identification of the seven maps}

We now pass from the coordinate algebra and its local geometry to the maps
themselves.  The analytic calculation supplies a uniform error bound; exact
canonical-ring linear algebra then reconstructs the ratios, and certified
branch cycles identify their moduli points.

\begin{computation}[Uniform truncation-error certificate]
\label{comp:hurwitz-tail}
The $47$ chart centres cover both half-triangles with pseudohyperbolic radius
less than $0.471$: an Arb calculation \cite{JohanssonArb} on geodesically
convex mesh cells in
the Klein disk model gives the upper bound $0.4708683715$.  At Cauchy radius $\rho=0.72$, the
modes $0,\ldots,60$, augmented by four normalization rows, have certified
minimum singular value at least $6.60\cdot10^{-4}$ in all seven classes.
Let $A$ be the augmented low-mode matrix, $h$ the norm bound for the
high-mode operator, and $\ell$ the bound for the low-mode transition
operator.  Parseval gives
$h\leq2.35\cdot10^{-11}$ and $\ell\leq1.38\cdot10^3$, so the explicitly
defined Schur-complement margin
\[
 m=\sigma_{\min}(A)-\frac{h\ell}{1-h}
\]
is positive in every class.

A bootstrap using Cauchy estimates at outer radius $0.99$ bounds the exact
normalized forms.
The final $N=700$, $Q=1280$ Acb calculation then evaluates all $1280$
discrete Fourier transform (DFT) output modes, not only those through the
polynomial cutoff.  Combining its
residual balls with the certified inverse and branch-jet conditioning gives
at least $75$ certified decimal digits in each of the first twenty
branch-normalized Taylor rows used to recover the canonical quadric and Petri
cubic.
\end{computation}

The exact equations are too large to reproduce usefully in print.  They are
stored in power or integral bases in
\path{hurwitz_canonical_models_candidate.json} and
\path{hurwitz_marked_points_candidate.json}.  Their importance here is that
they allow the covering function to be recovered by linear algebra internal
to the canonical ring, with no search for a low-degree plane ansatz.

\begin{proposition}[Exact degree-$23$ ratios]
\label{prop:hurwitz-maps}
Let $C\to\operatorname{Spec}A_{\mathrm H}$ be the reconstructed family of
canonical complete intersections, and let
$b,c:\operatorname{Spec}A_{\mathrm H}\to C$ be the two sections marking
the points with order-$23$ inertia.
There are explicit quintic canonical sections $N,D$ such that
\[
 u=\frac ND,
 \qquad \operatorname{div}(u)=23b-23c.
\]
In particular $u$ has degree $23$ on every geometric fibre.  Restriction to
the degree-$6$ factor of $A_{\mathrm H}$ and its six $K_0$-embeddings gives
six explicit conjugate maps.
\end{proposition}

\begin{proof}
The degree-$5$ piece of the canonical ring has dimension $27$.  The first
$23$ local coefficients at either marked point have rank $23$, and hence
\[
 B=H^0(5K_C-23b),\qquad
 C'=H^0(5K_C-23c)
\]
both have dimension $4$.  The degree-$10$ piece has dimension $57$.
Multiplication there turns the identities
\[
 B_i(MC')_j-B_j(MC')_i=0\qquad(0\leq i<j\leq3)
\]
into a linear system with $16$ unknowns.  Its rank is $15$; its kernel is
generated by an invertible matrix $M$.  The deterministic first row gives
$N=B_0$ and $D=(MC')_0$.  Exact local expansions give
\[
 \ord_b(N)=23,\quad D(b)\ne0,\qquad
 \ord_c(D)=23,\quad N(c)\ne0.
\]
Finally, saturating the ideal generated by the four sections of $B$ away
from $b$ gives the unit ideal, and the same calculation for $MC'$ away from
$c$ again gives the unit ideal.  There is therefore no residual common base
point, zero, or pole.  This proves the divisor identity.
\end{proof}

It remains to put the third branch point at $1$.  The raw value is
$\lambda=u(a)\in L$, where $a$ is any point above the order-$2$ branch
point.  Its coordinates in the chosen power basis have a large common
denominator, so ordinary coordinatewise recognition is ineffective.

\begin{computation}[Third branch normalization]
\label{comp:third-branch}
At $71$ rational primes splitting completely in $L$, the unique nonzero
linear factor of multiplicity $8$ in the reduced branch discriminant is
computed in all twelve residue embeddings and interpolated into the power
basis.  Their CRT modulus has $404$ decimal digits.  A $13$-dimensional LLL
calculation \cite{LLL} gives a $1129$-bit common-denominator vector; the next lattice
vectors have at least $1246$ bits.  This value agrees with the independent
$1024$-bit Arb computation to $1.7\cdot10^{-101}$ in integral-basis
coordinates.  It also agrees exactly at $24$ further completely split primes
that were withheld from CRT and LLL.

\smallskip
For a characteristic-zero check, project the canonical complete intersection
birationally to a plane sextic.  After cancellation, $N$ and $D$ become plane
forms of degree $7$.  The resultants of the plane sextic with a generic fibre
and with $N-\lambda D$ both have degree $42$.  Their gcds with their
$x$-derivatives have degrees $6$ and $14$, respectively; the common factor has
degree $6$, and the additional degree-$8$ quotient is squarefree.  The only
other possible source of such a repeated projected root is two distinct plane
points with the same $x$-coordinate.  Magma excludes this after good reduction
at the completely split prime $863153$: in residue embeddings
$2,3,4,5,9,12$, every irreducible factor of the degree-$8$ quotient gives a
degree-$1$ gcd in the remaining coordinate.  Since the resultant, gcd, and
squarefree-factor degrees are preserved, a characteristic-zero common divisor
of degree at least $2$ would remain one after reduction.  Thus the third fibre
has eight distinct double points and seven simple points.  Independently, the
projective critical ideal has length $15=7+8$ in all twelve residue embeddings.

\smallskip
With this value, the six maps are normalized by
\[
 \beta=\frac{N}{\lambda D};
\]
their branch values are $0,1,\infty$, with passport
\[
 (23),\qquad (2^8 1^7),\qquad (23).
\]
\end{computation}

The exact field element $\lambda$, all reconstruction and holdout residues,
and the sparse sections $N,D,\lambda D$ are stored in
\path{hurwitz_degree23_branch_candidate.json} and
\path{hurwitz_degree23_maps_candidate.json}.  The degree and the two totally
ramified fibres in \cref{prop:hurwitz-maps} are characteristic-zero exact
statements, as is the third-fibre multiplicity in
\cref{comp:third-branch}.  The uniform Taylor-tail estimate is supplied by
\cref{comp:hurwitz-tail}.  The filenames retain the word \emph{candidate} to
record that LLL and bounded-denominator reconstruction were the discovery
route for the displayed coefficients; the following independent cover
identification removes that reconstruction assumption from the final
Hurwitz-algebra statement.

\smallskip
\begin{computation}[Exact eliminant and certified branch cycles]
\label{comp:hurwitz-branch-cycles}
In the affine canonical chart $X_0=1$, eliminate $X_3=z$ from the quadric and
Petri cubic.  If their linear subresultant is $L_1(x,y)z+L_0(x,y)$, then every
serialized quintic is linear in $z$, so a section $S_0+zS_1$ becomes the
degree-$7$ plane form $S_0L_1-S_1L_0$.  Eliminating $y$ from the plane sextic
and $N-t\lambda D$ gives a raw resultant $R_t(x)$ of degree $42$.  The exact
gcd of $R_2$ and $R_3$ has degree $19$ and divides every $R_t$; after division,
the values at $t=2,\ldots,8$ interpolate
\[
 P(t,x)\in L[t,x],\qquad \deg_tP=6,\quad \deg_xP=23.
\]
The unused value $t=9$ agrees exactly.  At $t=0,1,\infty$, the gcds with the
$x$-derivative have degrees $22,8,22$, and the degree-$8$ gcd is squarefree.

\smallskip
For each of the six embeddings of $L$ over the chosen embedding of $K_0$, Arb
isolates the $23$ roots over $t_*=1/2+2i$ and continues them around rational
polygonal loops about $0,1,\infty$.  On every straight path segment, each
sheet is enclosed in a disk of radius $R$ with fixed inverse slope $Y$.
Writing
\[
 \eta=\sup|Y P(t,c)|,
 \qquad q=\sup|1-Y P_x(t,x)|,
\]
the segment is accepted only when $q<1$ and $\eta+qR<R$, all $23$ disks are
pairwise disjoint, and the endpoint root balls lie in their unique disks.
Otherwise the segment is bisected.  Thus the continuation is certified by a
uniform contraction, not by nearest-neighbour matching.  The map on the
degree-one component (Nielsen class $6$) is checked independently from its
integral minimal-degree equation by the same method.

\smallskip
The full run contains $150145$ accepted tubes and has maximum dyadic depth
$18$.  In the sextic-embedding order the Nielsen IDs are
\[
 (7,4,1,5,3,2).
\]
Together with the independently checked degree-one map, whose Nielsen ID is
$6$, the seven triples
have cycle shapes
\[
 (23),\qquad(1^7 2^8),\qquad(23),
\]
their products are the identity, and every pair generates a permutation
group of order $10200960$.  After conjugating the order-$23$ generator to the
fixed $y$, centralizer translation matches the seven involutions with Nielsen
representatives $1,\ldots,7$, exactly once each.  The independent normalization
of the degree-one map orders the two order-$23$ endpoints oppositely, so its
comparison
uses the inverse of its loop about $0$.
\end{computation}

\begin{corollary}[Identification of the finite Hurwitz scheme]
\label{cor:hurwitz-identification}
The finite étale scheme
$\mathcal H_J=\operatorname{Spec}(K_0\times L)$ is isomorphic to
$\mathcal H_{\mathrm{in}}$.  Thus $A_{\mathrm H}=K_0\times L$ is the
coordinate algebra of the seven-point inner Hurwitz scheme of type
$(2A,23A,23B)$.  In particular, all seven reconstructed maps have
monodromy $M_{23}$.
\end{corollary}

\begin{proof}
The group-side Nielsen calculation gives exactly seven inner classes.
\Cref{comp:hurwitz-branch-cycles} produces one exact cover in each of them.
They are therefore distinct and exhaustive.  The induced map between the
two reduced finite étale $K_0$-schemes is a bijection on geometric points,
hence an isomorphism.  This identifies the coordinate algebra without an
a priori height bound for the preceding LLL discovery.
\end{proof}

\begin{proof}[Proof of \cref{thm:hurwitz-summary}]
The coordinate algebra and smooth canonical curves are given by
\cref{comp:hurwitz-coordinate-algebra}.  The exact maps are constructed in
\cref{prop:hurwitz-maps,comp:third-branch}; their certified branch cycles
give the identification with $\mathcal H_{\mathrm{in}}$ in
\cref{comp:hurwitz-branch-cycles,cor:hurwitz-identification}.  Finally,
\cref{prop:hurwitz-galois-closure} proves the $S_6$ assertion.
\end{proof}

\subsection{Descent of finite invariants}

We first specify what \emph{descent} means here.  For a finite set $S$, let
$\underline S_{K_0}=\coprod_{s\in S}\operatorname{Spec}K_0$ be the
associated constant finite étale scheme.  An assignment
\[
 f:\mathcal H_{\mathrm{in}}(\overline K_0)\longrightarrow S
\]
descends over $K_0$ if it is invariant under
$\operatorname{Gal}(\overline K_0/K_0)$, equivalently if it is induced by a
$K_0$-morphism
$\mathcal H_{\mathrm{in}}\to\underline S_{K_0}$.  When
$S=\{0,1\}$, such morphisms are equivalent to idempotents of the coordinate
algebra $A_{\mathrm H}$.

We next define the finite-group assignments to be tested.  Let $G=M_{23}$,
let $(x,y,z)$ be one of the seven normalized triples, and put
$Y=\langle y\rangle$ and $Z=\langle z\rangle$.  The two-endpoint order
count is
\[
 \nu(Y,Z)=\#\bigl\{(a,b)\in(\FF_{23}^{\times})^2:
                 \ord(y^az^b)=23\bigr\}.
\]
For $c\in\{y,z\}$ and $d\in\FF_{23}^{\times}$, put
\[
 a_c(d)=\ord\bigl(xx^{c^d}\bigr),\qquad
 \kappa_m=\sum_{c\in\{y,z\}}
   \#\{d\in\FF_{23}^{\times}:a_c(d)=m\}\quad(m=2,4).
\]
Changing a generator of $Y$ or $Z$ only reindexes the exponents, and the
sum over $c\in\{y,z\}$ is unchanged when the two endpoints are exchanged.

\smallskip
Write $N(Y)=N_G(Y)$, use the convention $x^g=g^{-1}xg$, and set
\[
 \Omega=x^{N(Y)}\cap x^{N(Z)}.
\]
The stabilizer of $x$ in either normalizer is trivial: its order divides
both $|N_G(Y)|=253$ and $|C_G(x)|=2688$.  Thus, for $w\in\Omega$, the
transporter sets from $x$ to $w$ in $N(Y)$ and $N(Z)$ are singletons; denote
their elements by $n_Y(w)$ and $n_Z(w)$.  The relative transporter
\[
 c_w=n_Y(w)n_Z(w)^{-1}
\]
centralizes $x$.  Let $A_x$ be the set of eight two-element orbits of
$\langle x\rangle$ on the moved letters.  The induced permutation
$\overline c_w$ lies in
\[
 C_G(x)/\langle x\rangle\cong 2^3{:}\PSL_3(2)
       \cong\operatorname{AGL}(3,2),
\]
acting faithfully on $A_x$.  Define
\[
 \mu=\#\{\overline c_w:w\in\Omega\},\qquad
 \Theta=\sum_{w\in\Omega}[\overline c_w]
       \in\FF_2[C_G(x)/\langle x\rangle].
\]
Finally, let $\epsilon$ be the standard augmentation homomorphism of this
group algebra.  Then $\epsilon(\Theta)=|\Omega|\bmod2$.  The Hamming weight
$\operatorname{wt}(\Theta)$ is the number of nonzero coefficients of
$\Theta$ in the group basis.

\smallskip
\begin{computation}[Descent of the relative-transporter invariants]
\label{comp:relative-transporter-descent}
In the sextic embedding order $(7,4,1,5,3,2)$, the exact values of the
two-endpoint order count $\nu$, the affine-image count $\mu$, and the Hamming
weight of $\Theta$ are respectively
\[
\begin{split}
 \nu|_L&=(42,46,54,46,28,54),\\
 \mu|_L&=(14,28,16,28,31,16),\\
 \operatorname{wt}(\Theta)|_L&=(11,27,15,27,25,15).
\end{split}
\]
The cyclic-conjugacy order counts $(\kappa_2,\kappa_4)$ are likewise
\[
 ((4,12),(2,6),(0,16),(2,6),(0,8),(0,16)).
\]
None is constant.  By contrast,
\[
 \epsilon(\Theta)|_L=(1,1,1,1,1,1),
 \qquad \epsilon(\Theta)|_{K_0}=0.
\]
Moreover, on all seven geometric points,
\[
 \boldsymbol1_{\{\nu=32\}}
 =\boldsymbol1_{\{\mu=17\}}
 =\boldsymbol1_{\{(\kappa_2,\kappa_4)=(0,12)\}}
 =\boldsymbol1_{\{\epsilon(\Theta)=0\}}.
\]
\end{computation}

\smallskip
\begin{proposition}[Component idempotents detected by finite invariants]
\label{prop:relative-transporter-descent}
Let $e_1=(1,0)$ and $e_6=(0,1)$ be the primitive idempotents of
$A_{\mathrm H}=K_0\times L$.  The four Boolean assignments in
\cref{comp:relative-transporter-descent} descend to morphisms to
$\underline{\{0,1\}}_{K_0}$ and all correspond to
$e_1$; the augmentation corresponds to $e_6=1-e_1$.

The assignments $\nu$, $\mu$, and $(\kappa_2,\kappa_4)$ do not descend on
$\operatorname{Spec}L$ to morphisms into the corresponding constant finite
schemes.  Moreover $\Theta$ cannot descend as a section of a lisse
$\FF_2$-sheaf whose transition maps permute the natural group-algebra basis.
\end{proposition}

\begin{proof}
Because $L$ is a field, $\operatorname{Spec}L$ is connected.  Every morphism
from it to a constant finite étale scheme is therefore constant, whereas the
displayed values of the three assignments vary among its six geometric
points.  A permutation of the group-algebra basis preserves Hamming weight,
so the varying weights also rule out the stated basis-compatible descent of
$\Theta$.  The four equality predicates are $1$ only at ID $6$, while the
augmentation is $1$ precisely on the other six points.  Under the
correspondence between Boolean morphisms and idempotents, they are therefore
$e_1$ and $e_6$, respectively.
\end{proof}

This interpretation uses the already certified $1+6$ decomposition; the
branch-cycle calculation remains the independent identification of the
Hurwitz scheme.  A conceptual explanation would have to construct
$\epsilon(\Theta)$ in a Galois-compatible characteristic-$23$ model and prove
directly that the resulting Boolean morphism cuts out
$\operatorname{Spec}L\subset\mathcal H_{\mathrm{in}}$, while discarding the
richer assignments that fail to descend.

\smallskip
As a separate global application of the minimal equation, we now record a
uniform specialization theorem.  It does not enter the preceding analysis of
the fixed Nielsen class.

\section{Arithmetic specializations of the minimal equation}
\label{sec:unramified-specializations}

At $T=0$ the specialized polynomial is necessarily reducible and
non-squarefree: it records the decomposition and inertia groups at a branch
point, not the generic monodromy.  Away from the branch locus the same
degree-$23$ equation returns to the full Mathieu group.

We first record the specialization fact that identifies the ambient group.

\begin{lemma}[Specialization subgroup]\label{lem:specialization-subgroup}
Let $Q(T,x)\in\QQ[T,x]$ have splitting field $L/\QQ(T)$ with Galois group
$G$.  If the leading coefficient and discriminant in $x$ are nonzero at
$t\in\QQ$, then the splitting field of $Q(t,x)$ has Galois group isomorphic,
in its action on the roots, to a subgroup of $G$.
\end{lemma}

\begin{proof}
Work over $R=\QQ[T]_{(T-t)}$.  After dividing by its unit leading
coefficient, $Q$ is monic with unit discriminant.  The normalization in $L$
is unramified above $T=t$.  The residue extension at a chosen prime has
Galois group its decomposition group, a subgroup of $G$.  The distinct
specialized roots generate that residue extension, proving the claim.
\end{proof}

For an integer $t\equiv2\pmod{319}$, let $c_t$ be the positive gcd of the
coefficients of $P(t,x)$ and set
\[
 f_t(x)=c_t^{-1}P(t,x)\in\ZZ[x].
\]

\begin{proof}[Proof of Theorem~\ref{thm:specialization}]
Write $t=2+319n$.  Since $319=11\cdot29$,
\[
 P(t,x)\equiv P(2,x)\pmod{11},\qquad
 P(t,x)\equiv P(2,x)\pmod{29}.
\]
The leading coefficient is $5$ modulo $11$ and $22$ modulo $29$, so neither
prime divides $c_t$ and $f_t$ retains degree $23$.

Modulo $29$, the monic coefficient vector in increasing order is
\[
\begin{split}
(5,3,7,15,21,21,17,9,1,28,21,25,{}&3,14,20,28,\\
 &5,27,24,16,9,20,18,1).
\end{split}
\]
Exact powering gives
\[
 x^{29^{23}}\equiv x\pmod{\bar f_{t,29}},\qquad
 \gcd(\bar f_{t,29},x^{29}-x)=1.
\]
Rabin's criterion proves irreducibility because $23$ is prime.  Thus $f_t$
is irreducible over $\QQ$, and Frobenius supplies a $23$-cycle.

Modulo $11$ one has
\[
 \bar f_{t,11}=g_2g_7g_{14},
\]
where
\begin{align*}
g_2={}&x^2+8x+3,\\
g_7={}&x^7+6x^6+7x^5+6x^4+8x^2+3x+1,\\
g_{14}={}&x^{14}+8x^{13}+6x^{12}+8x^{11}+6x^{10}+6x^9
 +3x^8+2x^7\\
&\quad+10x^6+3x^5+x^4+8x^3+10x^2+7x+6.
\end{align*}
The factors are distinct and irreducible, so Frobenius supplies an element
of cycle type $(2,7,14)$ and order $14$.  Squarefreeness also proves the
characteristic-zero discriminant is nonzero, and
\cref{lem:specialization-subgroup,thm:main} embeds
$H_t=\Gal(f_t/\QQ)$ in $M_{23}$.

The ATLAS maximal-subgroup orders are \cite{ATLAS,ATLASM23}
\[
 443520,\ 40320,\ 40320,\ 20160,\ 7920,\ 5760,\ 253.
\]
Only $253=23\cdot11$ is divisible by $23$, and it is odd.  A subgroup
containing both a $23$-cycle and an element of order $14$ is therefore not
proper.  Hence $H_t=M_{23}$.
\end{proof}

\begin{corollary}\label{cor:explicit}
The explicit primitive polynomial $f_2(x)=P(2,x)$ has Galois group $M_{23}$
over $\QQ$.
\end{corollary}

\subsection{Controlled ramification inside the progression}

Call a family mutually independent if every finite subfamily is linearly
disjoint over $\QQ$.

Huang et al. already deduce the abstract existence of an infinite mutually
independent family from their regular realization
\cite[Corollary~1.4(d)]{HJLPPZ}.  The refinement below chooses such a family
inside the explicit progression of \cref{thm:specialization} while forcing a
new, controlled inertia prime at each step.

\begin{theorem}\label{thm:distinct}
Among the splitting fields of $f_t$ with $t\equiv2\pmod{319}$ there is an
infinite mutually independent collection of $M_{23}$-extensions.  More
precisely, outside a finite set of primes, if $q\nmid319$ and
\[
 t\equiv2\pmod{319},\qquad v_q(t)=1,
\]
then inertia at $q$ is generated by an element of class $2A$.
\end{theorem}

\begin{proof}
Beckmann's specialization inertia theorem
\cite[Proposition~4.2]{Beckmann}, applied at $T=0$, says that away from
finitely many primes, meeting the branch point with multiplicity $m$
produces inertia conjugate to $\langle g^m\rangle$ for local monodromy $g$.
Enlarge this exceptional set to contain $23$.  If $v_q(t)=1$, then
$T^2+23$ is a unit at the specialization, and the inertia group is the
class-$2A$ group $\langle g\rangle$.

The Chinese remainder theorem gives
\[
 t\equiv2\pmod{319},\qquad t\equiv q\pmod{q^2}.
\]
Choose each new $q$ outside the bad set and the ramified primes of all
previous fields.  The new field ramifies at $q$ while the preceding fields
do not, so the fields are pairwise distinct.

For mutual independence, suppose $L_1,\ldots,L_r$ are linearly disjoint and
$K_r=L_1\cdots L_r$, with Galois group $M_{23}^r$.  For a new field $L$,
$K_r\cap L$ is Galois over $\QQ$.  Simplicity of $M_{23}$ makes it either
$\QQ$ or $L$.  In the latter case a surjection
$M_{23}^r\twoheadrightarrow M_{23}$ must factor through one coordinate:
the coordinate images are commuting normal subgroups, and simplicity with
trivial center permits only one nontrivial image.  Then $L$ equals a previous
$L_i$, contradicting distinctness.  Hence $K_r\cap L=\QQ$.
\end{proof}

\smallskip
\begin{computation}[A second certified field]\label{comp:t3830}
The integer
\[
 3830\equiv2\pmod{319},\qquad 3830\equiv5\pmod{25}
\]
lies in the progression and meets $T=0$ $5$-adically with multiplicity one.
Over $\FF_5$, the repeated part
\[
 \gcd(F(0,V),F_V(0,V))
\]
has degree $8$, is squarefree, and is coprime to $F_T(0,V)$.  Exact
maximal-order computations give
\[
 v_5(\operatorname{Disc}(\QQ[x]/(f_2)))=0,
 \qquad
 v_5(\operatorname{Disc}(\QQ[x]/(f_{3830})))=8.
\]
Thus the $t=3830$ splitting field ramifies at $5$ while the $t=2$ field does
not.  The natural degree-$23$ action is faithful, so the same conclusion
holds for their Galois closures: inertia in the closure of the unramified
root field fixes every point, hence lies in the core of a point stabilizer,
which is trivial.
\end{computation}

The local analysis and specialization theorem answer complementary
questions.  The branch fibre reveals the centralizer and residue arithmetic
after a $2A$ collision; the progression gives smooth arithmetic fibres with
the full generic group.

\section{Open questions}
\label{sec:open}

The calculations leave two concrete problems of different kinds.  The first
asks whether the branch-fibre tower is arithmetically forced by its local
data.  The second asks for the missing conceptual explanation of the
degree-one Hurwitz component.
\begin{enumerate}
\item Classify the $2^4{:}\PSL_3(2)$-extensions of $\QQ$ unramified outside
  $\{2,3,23\}$ with the local groups and constituent fields of
  \cref{thm:ramification}; in particular, decide whether $M$ is unique.
\item Construct the augmentation $\epsilon(\Theta)$ intrinsically and in a
  Galois-compatible way from characteristic-$23$ models of the seven covers,
  and prove that its Boolean morphism cuts out the degree-six component
  without first using the exact branch cycles.  The richer two-endpoint and
  cyclic-conjugacy data cannot themselves descend, by
  \cref{prop:relative-transporter-descent}.  More concretely, resolve the
  $E_8$ and $A_2+A_6$ deformations of
  \cref{comp:hurwitz-pointed-23} as pointed stable $M_{23}$-maps and identify
  the resulting binary cycle with the local residue-degree parity
  $f(P)-1$.
\end{enumerate}

\clearpage
\section{Exact data and certificates}
\label{sec:certificates}

The repository accompanying this paper is
\begin{center}
\url{https://github.com/research-reflexivity/frontier-math-m23-cover-investigation}.
\end{center}
It contains canonical JSON tables for $F$, $P$, and the pencil
$(J_0,J_1)$.  The convention for the minimal-degree equation is
\[
 \texttt{coefficients\_T\_then\_W[i][j]}=[T^iW^j]P.
\]
\par\noindent
\begin{minipage}{\textwidth}
The polynomial has $120$ nonzero coefficients and content $1$; its
coefficient of $T^4W^{23}$ is
\par\vspace{4pt}
\centering
$838610528294018704285119511337682389342017079161221371107109332287875$.
\end{minipage}
\par\vspace{4pt}
The principal input fingerprints are
\par\vspace{4pt}
\begingroup
\centering\small
\path{data/Fint_coefficients_Z.json}\par
\texttt{27dd1fd0de4f1c350f5a07ead6d5b747d7a8f34e4146ec42fa953f1536a65102}\par\vspace{3pt}
\path{data/optimal_23_4_Z.json}\par
\texttt{cfb4185b2800744d524d87b3f08f5459cd38d42afb3787c11fa1e2aa3f25ee84}\par\vspace{3pt}
\path{data/optimal_degree4_pencil.json}\par
\texttt{fdf571e4f47effc527e1726bdd99472be2c5e3390e45133fb3d7a3aab592619d}\par\vspace{3pt}
\path{data/canonical_quadric_Q.json}\par
\texttt{352b978369353a1cb57a1a3f54edd6b7c024896edf78c38c1601a06f395ac943}\par\vspace{3pt}
\path{data/hurwitz_degree23_branch_candidate.json}\par
\texttt{cfbce099e3fa867b7c5ba29aa73f02f11d386cf60b00a506d1166d4c65639f62}\par\vspace{3pt}
\path{data/hurwitz_degree23_maps_candidate.json}\par
\texttt{86cdbd314d01a9123aed0b6890fa4316b4303782b1d69563e020b6ede62a6c2a}\par
\endgroup
\par\vspace{4pt}

The single command
\par\vspace{4pt}
{\centering\texttt{make verify-all}\par}
\par\vspace{4pt}
runs the complete open-source suite.  The optional target
\texttt{make verify-magma} repeats the minimal-model identity,
irreducibility, canonical-quadric, and unpointed $23$-inertia obstruction checks
in Magma, together with the twelve-embedding Hurwitz critical-fibre and
sextic Galois-closure checks.
The recorded environment uses SageMath 10.9 \cite{SageMath}, GAP 4.14.0
\cite{GAP}, Singular 4.4.1 \cite{Singular}, and PARI/GP 2.17.4
\cite{PARI}; the independent Magma records use version 2.29-9
\cite{Magma}.
The opt-in target \texttt{make verify-hurwitz-third-fiber-exact} recomputes
the longer characteristic-zero resultant certificate.  The main targets are:

{\small
\begin{longtable}{@{}
  >{\raggedright\arraybackslash}p{.27\textwidth}
  >{\raggedright\arraybackslash}p{.48\textwidth}
  >{\raggedright\arraybackslash}p{.15\textwidth}@{}}
\toprule
Result & Certificate & Engine\\
\midrule
\endfirsthead
\toprule
Result & Certificate & Engine\\
\midrule
\endhead
Newton-polygon valuations, nodes, adjoint pencil &
\path{verify_boundary_valuations.gp}, \path{verify_nodes_mod31.sing},
\path{verify_adjoint_mod31.py} & PARI, Singular, Python\\
Minimal equation and function-field identity &
\path{verify_optimal_23_4.py}, \path{verify_sage.py},
\path{verify_optimal_23_4.m} & Singular, SageMath, Magma\\
Specialization progression and ramification &
\path{verify_progression_319.*}, \path{verify_ramified_t3830.*} &
Python, PARI\\
Transport of the fibre over $T=0$ &
\path{verify_special_fibre_bridge.py} & SageMath\\
Branch cycles and Fano incidence &
\path{certify_branch_cycle_description.*},
\path{certify_special_fibre_fano_*} & GAP, PARI\\
Affine field tower and final quadratic layer &
\path{certify_special_fibre_field_tower.*},
\path{certify_vanishing_cycle_class.*} & GAP, PARI\\
Ramification, flags, and complex-conjugation comparison &
\path{certify_special_fibre_ramification.*},
\path{certify_local_flag_places.*},
\path{certify_real_frame_cocycle.*} & GAP, PARI\\
Geometric reflection obstruction and complex-conjugation comparison &
\path{certify_geometric_reflection_*} & GAP, Magma\\
Canonical quadric and ruling field &
\path{verify_canonical_quadric.py},
\path{verify_canonical_quadric.m} & SageMath, Magma\\
Unpointed $23$-inertia obstruction &
\path{certify_gauss_prolongation_obstruction.g},
\path{certify_gauss_prolongation_obstruction.m} & GAP, Magma\\
Coordinate algebra of the Hurwitz scheme, canonical curves, and marked points &
\path{verify_hurwitz_*_candidate.py} & SageMath\\
Uniform atlas cover and truncation-error bounds &
\path{verify_hurwitz_tail_geometry.py},
\path{verify_hurwitz_tail_summary.py} & Arb, Acb, Python\\
Degree-$23$ ratios, third branch value, and complete passport &
\path{verify_hurwitz_degree23_*.py},
\path{verify_hurwitz_degree23_third_fiber.sage},
\path{verify_hurwitz_degree23_branch.m},
\path{verify_hurwitz_degree23_geometry.m} & SageMath, Magma\\
Exact degree-$23$ eliminant and all seven branch-cycle triples &
\path{verify_hurwitz_monodromy_eliminant_candidate.py},
\path{verify_hurwitz_branch_cycle_summary.py} & SageMath, Arb\\
Relative-transporter invariants and component-idempotent descent &
\path{certify_hurwitz_relative_transporter.g},
\path{verify_hurwitz_relative_transporter.py} & GAP, SageMath\\
Galois closure and six-point permutation action &
\path{verify_hurwitz_galois_closure.*} & SageMath, PARI, Magma\\
Local decomposition group and pointed reductions at $23$ &
\path{verify_hurwitz_local_23.py},
\path{verify_hurwitz_pointed_23.py} & SageMath\\
\bottomrule
\end{longtable}
}

Several certificate filenames predate the terminology adopted in this
paper and retain the string \texttt{special\_fibre}.  In those filenames it
refers to the characteristic-zero fibre over $T=0$, not to the closed fibre
of a model over a local base.

For the coordinate-transport certificate, SageMath reconstructs the fibres
over $T=0$ in both plane models,
checks the factor multiplicities and scalar identities, proves
$\gcd(J_1(0,V),H_7R_8)=1$, and computes the minimal polynomial of $J_0/J_1$
in each residue field.  Thus the certificate verifies the coordinate
transport coefficient by coefficient, including regularity and residue-field
generation.
The largest computations in the Fano suite are the direct degree-$1344$
splitting-field calculation for $R_8$, the comparison of the possible
septic/octic Frobenius types, the $9{,}889$ transvection primes below
$10^6$, and the square-class tests at fourteen primes splitting completely
in $L_0$.

\clearpage
\section{Scope, provenance, and disclosure}
\label{sec:scope}

The starting cover is not constructed in this paper.  The regular
$M_{23}$ realization, its source equation, branch cycles, and the
numerical computation of the seven complex covers and their Galois-fixed
Nielsen class are due to Huang, Jackson, Lee, Poonen, Pries, and Zhang
\cite{HJLPPZ}.  Their accompanying repository supplies the defining data
used here \cite{M23data}.

The contributions of the present paper fall into four groups.
\begin{enumerate}
\item We explicitly construct the degree-$4$ function $W$ whose existence
  and minimality were established in \cite[Remark~3.4]{HJLPPZ}, the
  bidegree-$(23,4)$ equation $P(T,W)$, and exact certificates that it defines
  the same function field.  We also identify the canonical ruling field that
  appears in their minimality argument.
\item We determine the arithmetic of the fibre over $T=0$: its Fano
  incidence, Gassmann pair, affine torsor, local field tower, ramification,
  and flags selected by complex conjugation and Frobenius.  We transport all
  of this structure exactly to the minimal-degree coordinate and prove the
  reflection and unpointed-inertia obstructions.
\item We reconstruct the coordinate algebra of the inner Hurwitz scheme,
  its seven canonical curves and degree-$23$ maps, their complete passports,
  and their branch cycles.  We determine the $S_6$ Galois action on the
  degree-six component, its local decomposition and pointed reductions at
  $23$, and the descent behavior of the relative-transporter invariants.
\item We prove the explicit uniform specialization theorem and controlled
  ramification inside its progression, and construct within that progression
  an infinite mutually independent family.  The abstract existence of such
  families is already \cite[Corollary~1.4(d)]{HJLPPZ}.
\end{enumerate}

We do not reprove that the fixed Nielsen class is rational; that is a
theorem of \cite{HJLPPZ}.  Nor do we give a criterion forcing the degree-one
component.  The local residue-degree parity agrees with the binary
augmentation, but we do not construct that agreement intrinsically from
pointed stable $M_{23}$-maps.  Resolving the $E_8$ and $A_2+A_6$
deformations in a common Mathieu frame remains open.

The third branch fibre of the six conjugate maps is verified exactly in
characteristic zero and independently in all twelve residue embeddings of
a completely split prime.  This distinction matters: numerical and lattice
reconstruction discover the coefficients, while the subsequent exact
elimination and branch-cycle certificates identify the covers.

The discovery and drafting process used substantial AI assistance.  Every
new computational assertion is covered by an exact certificate described in
\cref{sec:certificates}.

\subsection*{Acknowledgements}

This work analyzes the construction of Huang, Jackson, Lee, Poonen, Pries,
and Zhang and depends on their making the defining data publicly available
in the repository cited as \cite{M23data}.

\clearpage
\begin{thebibliography}{99}

\bibitem{HJLPPZ}
X.~Huang, B.~Jackson, K.-H.~Lee, B.~Poonen, R.~Pries, and S.~Zhang,
\emph{The Mathieu group $M_{23}$ is a Galois group over $\QQ$},
arXiv:2608.08538v1 (2026),
\url{https://arxiv.org/abs/2608.08538}.

\bibitem{M23data}
S.~Zhang et al., \emph{m23isgalois}, data repository,
\url{https://github.com/shaowuz/m23isgalois}.

\bibitem{BraidOrbits}
F.~H\"afner,
\emph{Braid orbits and the Mathieu group $M_{23}$ as Galois group},
arXiv:2202.08222 (2022).

\bibitem{KMSV}
M.~Klug, M.~Musty, S.~Schiavone, and J.~Voight,
\emph{Numerical calculation of three-point branched covers of the
projective line},
LMS J.~Comput. Math. \textbf{17} (2014), 379--430.

\bibitem{BelyiDB}
M.~Musty, S.~Schiavone, J.~Sijsling, and J.~Voight,
\emph{A database of Belyi maps},
Open Book Ser. \textbf{2} (2019), 375--392.

\bibitem{FriedVolklein}
M.~D.~Fried and H.~V\"olklein,
\emph{The inverse Galois problem and rational points on moduli spaces},
Math. Ann. \textbf{290} (1991), 771--800.

\bibitem{Rosenlicht}
M.~Rosenlicht,
\emph{Equivalence relations on algebraic curves},
Ann. of Math. (2) \textbf{56} (1952), 169--191.

\bibitem{Perlis}
R.~Perlis,
\emph{On the equation $\zeta_K(s)=\zeta_{K'}(s)$},
J. Number Theory \textbf{9} (1977), 342--360.

\bibitem{LLL}
A.~K.~Lenstra, H.~W.~Lenstra, Jr., and L.~Lov\'asz,
\emph{Factoring polynomials with rational coefficients},
Math. Ann. \textbf{261} (1982), 515--534.

\bibitem{JohanssonArb}
F.~Johansson,
\emph{Arb: Efficient arbitrary-precision midpoint-radius interval arithmetic},
IEEE Trans. Comput. \textbf{66} (2017), 1281--1292.

\bibitem{GLS}
G.-M.~Greuel, C.~Lossen, and E.~Shustin,
\emph{Introduction to singularities and deformations}, 2nd ed.,
Springer Monographs in Mathematics, Springer, Cham, 2025.

\bibitem{ATLAS}
J.~H.~Conway, R.~T.~Curtis, S.~P.~Norton, R.~A.~Parker, and R.~A.~Wilson,
\emph{Atlas of finite groups}, Clarendon Press, Oxford, 1985.

\bibitem{ATLASM23}
The ATLAS of Finite Group Representations,
\emph{Mathieu group $M_{23}$: maximal subgroups},
\url{https://brauer.maths.qmul.ac.uk/Atlas/v3/spor/M23/}.

\bibitem{Beckmann}
S.~Beckmann,
\emph{On extensions of number fields obtained by specializing branched
coverings},
J. Reine Angew. Math. \textbf{419} (1991), 27--54.

\bibitem{DebesGhazi}
P.~D\`ebes and N.~Ghazi,
\emph{Galois covers and the Hilbert--Grunwald property},
Ann. Inst. Fourier (Grenoble) \textbf{62} (2012), no.~3, 989--1013.

\bibitem{Legrand}
F.~Legrand,
\emph{Specialization results and ramification conditions},
Israel J. Math. \textbf{214} (2016), no.~2, 621--650.

\bibitem{LegrandEmbedding}
F.~Legrand,
\emph{On finite embedding problems with abelian kernels},
J. Algebra \textbf{594} (2022), 51--73.

\bibitem{SGA1}
A.~Grothendieck and M.~Raynaud,
\emph{Rev\^etements \'etales et groupe fondamental (SGA 1)},
Documents Math\'ematiques \textbf{3}, Soci\'et\'e Math\'ematique de France,
2003; arXiv:math/0206203.

\bibitem{Magma}
W.~Bosma, J.~Cannon, and C.~Playoust,
\emph{The Magma algebra system I: The user language},
J. Symbolic Comput. \textbf{24} (1997), 235--265.

\bibitem{SageMath}
The Sage Developers,
\emph{SageMath, the Sage mathematics software system}, version 10.9, 2026,
\url{https://www.sagemath.org}.

\bibitem{GAP}
The GAP Group,
\emph{GAP---Groups, Algorithms, and Programming}, version 4.14.0, 2024,
\url{https://www.gap-system.org}.

\bibitem{Singular}
W.~Decker, G.-M.~Greuel, G.~Pfister, and H.~Sch\"onemann,
\emph{Singular 4-4-1---A computer algebra system for polynomial computations},
2025, \url{https://www.singular.uni-kl.de}.

\bibitem{PARI}
The PARI Group,
\emph{PARI/GP}, version 2.17.4, Universit\'e de Bordeaux, 2026,
\url{https://pari.math.u-bordeaux.fr}.

\end{thebibliography}

\end{document}
